% =============================================================================
%  GARAH — MANUEL DE FORMATION UTILISATEUR
% =============================================================================
%
%  Compilation : pdfLaTeX (Overleaf : Menu > Compiler > pdfLaTeX).
%  Deux passes sont nécessaires pour la table des matières et les renvois.
%
%  ⚠️ AUCUNE IMAGE EXTERNE N'EST REQUISE POUR COMPILER.
%     Toutes les illustrations de ce manuel sont dessinées en TikZ. Elles sont
%     présentées comme des ILLUSTRATIONS PÉDAGOGIQUES : elles reproduisent
%     fidèlement la disposition des écrans réels, mais ce ne sont pas des
%     captures. Les emplacements où une vraie capture apporterait davantage
%     sont signalés par un cadre « CAPTURE À INSÉRER ».
%
% =============================================================================

\documentclass[11pt,a4paper]{report}

% =============================================================================
%  ⚠️  SI OVERLEAF DIT « YOUR COMPILE TIMED OUT » : LISEZ CECI.
% =============================================================================
%
%  Le plan gratuit d'Overleaf limite le temps de compilation. Les vingt
%  illustrations TikZ de ce manuel sont ce qui coûte le plus.
%
%  Passez la ligne ci-dessous à \illustrationsfalse : les figures sont alors
%  remplacées par des cadres, LEURS LÉGENDES RESTENT — elles décrivent ce que
%  la figure montrait — et le document compile en quelques secondes.
%
%  Vous obtenez un PDF complet, utilisable, où il ne manque que les dessins.
%  Remettez \illustrationstrue dès que vous compilez ailleurs : en local, sur
%  un plan payant, ou dans une chaîne d'intégration.
%
% =============================================================================
\newif\ifillustrations
\illustrationstrue          % <-- \illustrationsfalse si le temps manque

\usepackage{comment}
\ifillustrations
  \includecomment{illustration}
\else
  \excludecomment{illustration}
\fi


% --- Langue et encodage ------------------------------------------------------
\usepackage[utf8]{inputenc}
\usepackage[T1]{fontenc}
\usepackage[french]{babel}
\usepackage{lmodern}
% ⚠️ `microtype` retirée : elle affine la typographie, mais elle coûte du
%    temps de compilation — et le plan gratuit d'Overleaf en manque.
%    La remettre si le document compile chez vous avec de la marge.

% --- Mise en page ------------------------------------------------------------
\usepackage[a4paper,margin=2.4cm,headheight=15pt]{geometry}
\usepackage{fancyhdr}
\usepackage{titlesec}
\usepackage{enumitem}
\usepackage{booktabs}
\usepackage{array}
\usepackage{longtable}
\usepackage{amssymb}
\usepackage{pifont}
\usepackage{graphicx}
\usepackage{xcolor}

% --- Encadrés ----------------------------------------------------------------
% ⚠️ `breakable` SEULE, et non `most` : cette dernière charge une douzaine
%    de bibliothèques (raster, theorems, magazine…) dont aucune ne sert ici,
%    et chacune coûte du temps de compilation.
\usepackage[breakable]{tcolorbox}

% --- Dessins -----------------------------------------------------------------
\usepackage{tikz}
% Seules celles qui servent réellement : `positioning` pour les `right=of`,
% `arrows.meta` pour les pointes Stealth, `calc` pour la coordonnée calculée
% du cycle de commande.
\usetikzlibrary{positioning,arrows.meta,calc}

% --- Liens (toujours en dernier) --------------------------------------------
\usepackage[hidelinks,bookmarks=true]{hyperref}
\hypersetup{
  pdftitle={GARAH — Manuel de formation utilisateur},
  pdfsubject={Formation pratique à la plateforme GARAH},
  pdfkeywords={GARAH, formation, manuel, commerce, Douala, Bangui}
}

% =============================================================================
%  LA CHARTE GRAPHIQUE DE GARAH
%  Les valeurs viennent de charte/jetons.json — elles ne sont pas inventées.
% =============================================================================
\definecolor{garahMarque}{HTML}{31AEF3}   % le bleu de la marque
\definecolor{garahPrimaire}{HTML}{6366F1} % l'indigo des actions
\definecolor{garahAccent}{HTML}{AA3BFF}   % le violet du back-office
\definecolor{garahSlogan}{HTML}{047857}   % le vert du slogan
\definecolor{garahSucces}{HTML}{10B981}
\definecolor{garahAlerte}{HTML}{F59E0B}
\definecolor{garahDanger}{HTML}{EF4444}
\definecolor{garahEncre}{HTML}{0F172A}
\definecolor{garahAttenue}{HTML}{64748B}
\definecolor{garahFond}{HTML}{F8FAFC}
\definecolor{garahBordure}{HTML}{CBD5E1}
\definecolor{garahNuit}{HTML}{0C0C14}
\definecolor{garahSurfaceNuit}{HTML}{12121E}

% =============================================================================
%  TITRES
% =============================================================================
\titleformat{\chapter}[display]
  {\normalfont\huge\bfseries\color{garahEncre}}
  {\vspace*{-1.2cm}\normalsize\color{garahMarque}\MakeUppercase{\chaptertitlename}\ \thechapter}
  {8pt}{\Huge}
\titlespacing*{\chapter}{0pt}{0pt}{28pt}

\titleformat{\section}
  {\normalfont\Large\bfseries\color{garahEncre}}{\thesection}{0.8em}{}
\titleformat{\subsection}
  {\normalfont\large\bfseries\color{garahPrimaire}}{\thesubsection}{0.7em}{}

\pagestyle{fancy}
\fancyhf{}
\fancyhead[L]{\footnotesize\color{garahAttenue}GARAH — Manuel de formation}
\fancyhead[R]{\footnotesize\color{garahAttenue}\leftmark}
\fancyfoot[C]{\thepage}
\renewcommand{\headrulewidth}{0.4pt}

% =============================================================================
%  OUTILS PÉDAGOGIQUES
% =============================================================================

% --- Les pastilles numérotées ① ② ③ utilisées sur les illustrations ---------
\newcommand{\pastille}[2]{%
  {\setlength{\fboxsep}{1.7pt}%
   \colorbox{#1}{\color{white}\bfseries\scriptsize #2}}%
}
\newcommand{\num}[1]{\pastille{garahPrimaire}{#1}}
\newcommand{\numAlt}[1]{\pastille{garahAlerte}{#1}}

% --- Case à cocher pour les checklists --------------------------------------
\newcommand{\case}{$\square$~}

% --- Encadrés ---------------------------------------------------------------
\newtcolorbox{objectif}{
  breakable, colback=garahMarque!6, colframe=garahMarque,
  boxrule=0.9pt, arc=3pt, left=9pt, right=9pt, top=7pt, bottom=7pt,
  title={\bfseries Objectif de ce tutoriel}, fonttitle=\bfseries\small,
  coltitle=white, colbacktitle=garahMarque
}

\newtcolorbox{avant}{
  breakable, colback=garahFond, colframe=garahBordure,
  boxrule=0.8pt, arc=3pt, left=9pt, right=9pt, top=7pt, bottom=7pt,
  title={\bfseries Avant de commencer}, fonttitle=\bfseries\small,
  coltitle=garahEncre, colbacktitle=garahBordure!60
}

\newtcolorbox{attention}{
  breakable, colback=garahDanger!6, colframe=garahDanger,
  boxrule=0.9pt, arc=3pt, left=9pt, right=9pt, top=7pt, bottom=7pt,
  title={\bfseries \ding{55}~Attention}, fonttitle=\bfseries\small,
  coltitle=white, colbacktitle=garahDanger
}

\newtcolorbox{astuce}{
  breakable, colback=garahSucces!7, colframe=garahSucces,
  boxrule=0.9pt, arc=3pt, left=9pt, right=9pt, top=7pt, bottom=7pt,
  title={\bfseries \ding{43}~Astuce}, fonttitle=\bfseries\small,
  coltitle=white, colbacktitle=garahSucces
}

\newtcolorbox{retenir}{
  breakable, colback=garahAccent!6, colframe=garahAccent,
  boxrule=0.9pt, arc=3pt, left=9pt, right=9pt, top=7pt, bottom=7pt,
  title={\bfseries \ding{45}~À retenir}, fonttitle=\bfseries\small,
  coltitle=white, colbacktitle=garahAccent
}

\newtcolorbox{resultat}{
  breakable, colback=garahSucces!5, colframe=garahSucces!70,
  boxrule=0.9pt, arc=3pt, left=9pt, right=9pt, top=7pt, bottom=7pt,
  title={\bfseries \ding{51}~Résultat attendu}, fonttitle=\bfseries\small,
  coltitle=white, colbacktitle=garahSucces!70
}

\newtcolorbox{exercice}{
  breakable, colback=garahAlerte!7, colframe=garahAlerte,
  boxrule=0.9pt, arc=3pt, left=9pt, right=9pt, top=7pt, bottom=7pt,
  title={\bfseries Exercice pratique}, fonttitle=\bfseries\small,
  coltitle=white, colbacktitle=garahAlerte
}

\newtcolorbox{roleBox}{
  breakable, colback=garahPrimaire!5, colframe=garahPrimaire,
  boxrule=1pt, arc=4pt, left=11pt, right=11pt, top=9pt, bottom=9pt,
  title={\bfseries Votre rôle dans la plateforme}, fonttitle=\bfseries,
  coltitle=white, colbacktitle=garahPrimaire
}

\newtcolorbox{checklist}{
  breakable, colback=white, colframe=garahEncre,
  boxrule=1pt, arc=3pt, left=11pt, right=11pt, top=9pt, bottom=9pt,
  title={\bfseries Checklist de maîtrise}, fonttitle=\bfseries,
  coltitle=white, colbacktitle=garahEncre
}

% --- L'étape d'un tutoriel ---------------------------------------------------
\newcounter{etape}[section]
\newcommand{\etape}[1]{%
  \stepcounter{etape}%
  \vspace{9pt}%
  \noindent{\color{garahPrimaire}\rule{3pt}{13pt}}\hspace{6pt}%
  {\bfseries\large Étape \theetape\ --- #1}\par\vspace{4pt}%
}

% --- Le cadre « capture réelle à insérer » ----------------------------------
\newcommand{\captureAInserer}[2][3.4cm]{%
  \begin{center}
  \begin{tcolorbox}[width=\linewidth,height=#1,
    colback=garahFond,colframe=garahAttenue,boxrule=0.9pt,arc=4pt,
    halign=center,valign=center]
    {\small\bfseries\color{garahAttenue}CAPTURE D'ÉCRAN RÉELLE À INSÉRER}\\[4pt]
    {\footnotesize\color{garahAttenue}#2}
  \end{tcolorbox}
  \end{center}
}

% --- La légende d'une illustration pédagogique ------------------------------
\newcommand{\legendeIllus}[2]{%
  % Sans les dessins, la légende resterait suspendue dans le vide : on pose un
  % cadre à la place, pour que le renvoi « figure X » garde un objet.
  \ifillustrations\else
    \captureAInserer[2.6cm]{Figure~#1 --- illustration désactivée
      (voir l'interrupteur en tête de fichier). La légende ci-dessous décrit
      ce que la figure montre.}%
  \fi
  \vspace{2pt}
  \begin{center}
  \begin{minipage}{0.93\linewidth}
    \footnotesize\itshape
    \textbf{Figure~#1 --- } #2
    \par\vspace{2pt}
    \textcolor{garahAttenue}{\scriptsize Illustration pédagogique : reproduction
    fidèle de la disposition de l'écran, et non une capture.}
  \end{minipage}
  \end{center}
  \vspace{6pt}
}

% =============================================================================
\begin{document}
% =============================================================================

% -----------------------------------------------------------------------------
%  COUVERTURE
% -----------------------------------------------------------------------------
\begin{titlepage}
\centering
\vspace*{1.4cm}

% ⚠️ `scale` est indispensable : sans lui, les coordonnées ci-dessous sont
%    lues en centimètres et l'emblème ferait 12 cm de large — plus que la
%    moitié de la page.
\begin{tikzpicture}[scale=0.34]
  % L'emblème : le chariot de la marque, simplifié.
  \draw[garahMarque,line width=2.2pt,rounded corners=2pt]
    (0,0) -- (0,10) -- (2,10);
  \draw[garahMarque,line width=2.2pt] (2,10) -- (3.2,3.4);
  \draw[garahMarque,line width=2.2pt,rounded corners=1.4pt]
    (1.2,3.4) rectangle (12,7.2);
  \filldraw[garahMarque] (3.4,2) circle (1);
  \filldraw[garahMarque] (10,2) circle (1);
\end{tikzpicture}

\vspace{0.7cm}

{\fontsize{52}{56}\selectfont\bfseries\color{garahEncre}%
 G\textcolor{garahMarque}{AR}\textcolor{garahAccent}{AH}\par}

\vspace{4pt}
{\large\bfseries\color{garahSlogan}\MakeUppercase{Au-delà des frontières}\par}

\vspace{2.2cm}

{\Huge\bfseries\color{garahEncre}Manuel de formation\par}
\vspace{6pt}
{\Huge\bfseries\color{garahEncre}utilisateur\par}

\vspace{0.9cm}
\rule{0.55\linewidth}{1.4pt}
\vspace{0.7cm}

{\large\color{garahAttenue}
 Apprendre à utiliser la plateforme, écran par écran\par}

\vspace{2.4cm}

\begin{tcolorbox}[width=0.86\linewidth,colback=garahFond,
  colframe=garahBordure,arc=5pt,boxrule=1pt]
\centering\small
\textbf{Pour qui ?}\\[4pt]
Les clients \textbullet{} les responsables \textbullet{} les administrateurs\\
\textbf{et} les super administrateurs de GARAH
\\[9pt]
\textbf{Ce que vous saurez faire à la fin}\\[4pt]
Vous connecter, naviguer, effectuer vos tâches quotidiennes,\\
comprendre les résultats et vous dépanner seul.
\end{tcolorbox}

\vfill
{\footnotesize\color{garahAttenue}
 Corridor commercial Douala $\rightarrow$ Bangui\\
 Document de formation --- à lire l'application ouverte devant soi\par}
\end{titlepage}

% -----------------------------------------------------------------------------
\tableofcontents
\clearpage

% =============================================================================
\chapter*{Objectifs de la formation}
\addcontentsline{toc}{chapter}{Objectifs de la formation}
\markboth{Objectifs}{Objectifs}
% =============================================================================

Ce manuel n'explique pas comment GARAH est construit. Il explique
\textbf{comment s'en servir}.

À la fin de votre chapitre, vous devez être capable, seul et sans aide :

\begin{enumerate}[leftmargin=*,itemsep=4pt]
  \item de \textbf{comprendre votre rôle} : ce que la plateforme attend de
        vous, et ce qu'elle ne vous laissera pas faire ;
  \item de \textbf{vous connecter} et de retrouver votre chemin ;
  \item de \textbf{naviguer} dans les menus sans hésiter ;
  \item d'\textbf{effectuer vos tâches} quotidiennes de bout en bout ;
  \item de \textbf{comprendre ce que l'écran vous répond} ;
  \item de \textbf{vous dépanner} quand quelque chose ne marche pas ;
  \item de \textbf{devenir autonome}, c'est-à-dire de ne plus avoir besoin de
        ce manuel.
\end{enumerate}

\begin{retenir}
Ce manuel se lit \textbf{l'application ouverte devant soi}. Chaque tutoriel est
fait pour être exécuté en même temps qu'il est lu. Le lire sans rien faire
n'apprend presque rien.
\end{retenir}

% =============================================================================
\chapter*{Comment utiliser ce manuel}
\addcontentsline{toc}{chapter}{Comment utiliser ce manuel}
\markboth{Mode d'emploi}{Mode d'emploi}
% =============================================================================

\section*{Trouvez d'abord votre chapitre}

Vous n'avez pas besoin de tout lire. Repérez qui vous êtes :

\begin{center}
\begin{tabular}{@{}>{\bfseries}l l l@{}}
\toprule
Vous êtes\dots & Vous utilisez\dots & Allez au chapitre \\
\midrule
Un client        & la boutique ou l'application  & \ref{ch:client} \\
Un responsable   & le back-office               & \ref{ch:responsable} \\
Un administrateur& le back-office               & \ref{ch:admin} \\
Un super admin.  & le back-office               & \ref{ch:superadmin} \\
\bottomrule
\end{tabular}
\end{center}

\section*{Les repères visuels}

Ce manuel utilise toujours les mêmes signes. Apprenez-les une fois :

\begin{center}
\begin{tabular}{@{}c p{0.78\linewidth}@{}}
\toprule
\num{1} & \textbf{Une pastille bleue numérotée} désigne, sur une illustration,
          l'endroit exact où agir. Les numéros suivent l'ordre des étapes. \\[5pt]
\numAlt{!} & \textbf{Une pastille orange} signale une zone à ne pas modifier,
             ou une information à lire attentivement. \\[5pt]
\ding{55} & \textbf{Attention} : ce qui suit peut casser quelque chose, faire
            perdre des données, ou coûter de l'argent. \\[5pt]
\ding{43} & \textbf{Astuce} : un raccourci ou une habitude qui fait gagner du
            temps. \\[5pt]
\ding{45} & \textbf{À retenir} : le point essentiel du passage. \\[5pt]
\ding{51} & \textbf{Résultat attendu} : ce que vous devez voir à l'écran si
            tout s'est bien passé. \\
\bottomrule
\end{tabular}
\end{center}

\section*{À propos des illustrations}

\begin{attention}
Les figures de ce manuel sont des \textbf{illustrations pédagogiques}. Elles
reproduisent fidèlement la \emph{disposition} des écrans réels — l'ordre des
menus, la place des boutons, les libellés — mais ce ne sont pas des captures
d'écran.

Là où une véritable capture apporterait davantage, un cadre en pointillés
indique où l'insérer :

\captureAInserer[6em]{Exemple d'emplacement réservé.}

Un formateur qui prépare une session devrait remplacer ces cadres par les
captures de son propre environnement.
\end{attention}

% =============================================================================
\chapter{Découvrir GARAH}
\label{ch:decouvrir}
% =============================================================================

\section{Ce que fait la plateforme}

GARAH est une plateforme de commerce qui relie \textbf{Douala} (Cameroun) à
\textbf{Bangui} (République centrafricaine).

Le principe tient en une phrase :

\begin{center}
\begin{tcolorbox}[width=0.9\linewidth,colback=garahMarque!8,
  colframe=garahMarque,arc=4pt,boxrule=1pt]
\centering\large
Un client commande en ligne, paie par Mobile Money,\\
et \textbf{vient retirer sa marchandise} au point de récupération de son choix.
\end{tcolorbox}
\end{center}

Il n'y a \textbf{pas de livraison à domicile}. La marchandise voyage jusqu'à un
point de récupération, et le client s'y présente avec son \textbf{code de
retrait}.

\subsection{Le voyage d'une commande, en une image}

\begin{illustration}
\begin{center}
\begin{tikzpicture}[
  node distance=0pt,
  etape/.style={rectangle,rounded corners=3pt,draw=garahMarque,
                fill=garahMarque!8,line width=0.9pt,
                minimum height=13mm,text width=24mm,align=center,
                font=\footnotesize\bfseries},
  fleche/.style={-{Stealth[length=3mm]},garahAttenue,line width=1pt},
  acteur/.style={font=\scriptsize\itshape,text=garahAttenue}
]
\node[etape] (a) {Le client\\commande};
\node[etape,right=9mm of a] (b) {Il paie\\Mobile Money};
\node[etape,right=9mm of b] (c) {L'équipe\\prépare};
\node[etape,right=9mm of c] (d) {Le colis\\voyage};
\node[etape,right=9mm of d] (e) {Retrait\\avec le code};

\draw[fleche] (a) -- (b);
\draw[fleche] (b) -- (c);
\draw[fleche] (c) -- (d);
\draw[fleche] (d) -- (e);

\node[acteur,below=2mm of a] {Client};
\node[acteur,below=2mm of b] {Client};
\node[acteur,below=2mm of c] {Responsable};
\node[acteur,below=2mm of d] {Responsable};
\node[acteur,below=2mm of e] {Client + point};
\end{tikzpicture}
\end{center}
\end{illustration}
\legendeIllus{1.1}{Le parcours complet d'une commande. La ligne du bas indique
qui agit à chaque étape.}

\section{Les quatre acteurs}

La plateforme connaît exactement \textbf{quatre types de comptes}. Ils ne se
choisissent pas : ils sont attribués.

\begin{illustration}
\begin{center}
\begin{tikzpicture}[
  boite/.style={rectangle,rounded corners=4pt,line width=1pt,
                minimum height=26mm,text width=33mm,align=center,
                font=\footnotesize},
  titre/.style={font=\small\bfseries}
]
\node[boite,draw=garahMarque,fill=garahMarque!7] (cli) at (0,0)
  {\textbf{\normalsize Client}\\[4pt]
   Achète.\\[3pt]
   \scriptsize Boutique web\\ ou application Android};
\node[boite,draw=garahPrimaire,fill=garahPrimaire!7,right=6mm of cli] (res)
  {\textbf{\normalsize Responsable}\\[4pt]
   Fait le travail.\\[3pt]
   \scriptsize Back-office\\ (droits limités à son profil)};
\node[boite,draw=garahAccent,fill=garahAccent!7,right=6mm of res] (adm)
  {\textbf{\normalsize Administrateur}\\[4pt]
   Décide qui fait quoi.\\[3pt]
   \scriptsize Back-office\\ (équipe, profils, finance)};
\node[boite,draw=garahDanger,fill=garahDanger!7,right=6mm of adm] (sup)
  {\textbf{\normalsize Super admin.}\\[4pt]
   Surveille tout.\\[3pt]
   \scriptsize Back-office\\ + journal et sécurité};
\end{tikzpicture}
\end{center}
\end{illustration}
\legendeIllus{1.2}{Les quatre acteurs et l'interface de chacun.}

\begin{retenir}
Un compte \textbf{n'a pas de droits en propre} : il en reçoit par ses
\textbf{profils}. C'est la phrase affichée en haut de l'écran « Équipe » du
back-office :

\begin{center}
\small\ttfamily
ADMINISTRATEUR $\rightarrow$ PROFIL $\rightarrow$ RESPONSABLE\\
\rmfamily\footnotesize
décide qui fait quoi \quad un paquet de droits \quad fait le travail
\end{center}
\end{retenir}

\section{Les trois applications}

Une même plateforme, trois portes d'entrée :

\begin{center}
\begin{tabular}{@{}>{\bfseries}p{3.4cm} p{4.6cm} p{6.2cm}@{}}
\toprule
Application & Pour qui & Ce qu'on y fait \\
\midrule
La boutique web
  & Les clients, sur ordinateur ou téléphone
  & Parcourir, commander, suivre, discuter \\[4pt]
L'application Android
  & Les clients, sur téléphone
  & La même chose, en plus rapide et hors ligne pour le panier \\[4pt]
Le back-office
  & Responsables, administrateurs, super administrateurs
  & Tout le travail interne \\
\bottomrule
\end{tabular}
\end{center}

\begin{astuce}
La boutique et l'application ne sont pas la même chose enveloppée deux fois.
L'application garde votre panier même fermée, et ne redemande pas votre mot de
passe à chaque ouverture.
\end{astuce}

\section{Vocabulaire essentiel}

Ces mots reviennent partout. Apprenez-les maintenant, ils ne seront plus
réexpliqués.

\begin{longtable}{@{}>{\bfseries}p{3.6cm} p{11.2cm}@{}}
\toprule
Mot & Ce qu'il veut dire chez GARAH \\
\midrule
\endhead
Déclinaison
  & Une variante achetable d'un article : la même chaussure en blanc et en
    noir sont deux déclinaisons. C'est la déclinaison qui a un stock et un
    prix, pas l'article. \\[4pt]
Palier de prix
  & Le prix baisse par tranches de quantité : 1 à 6 pièces à un prix, 7 à 10 à
    un autre, 11 et plus à un troisième. \\[4pt]
Point de récupération
  & Le lieu où le client vient chercher sa marchandise. Il choisit le sien au
    moment de commander ; les frais d'acheminement en dépendent. \\[4pt]
Code de retrait
  & Le code unique remis au client. C'est le \textbf{seul} moyen de prouver que
    la marchandise a été remise. \\[4pt]
Expédition
  & Le voyage d'un ou plusieurs colis sur un itinéraire. \\[4pt]
Itinéraire
  & Le chemin suivi par les colis, avec ses points de transit. \\[4pt]
Conversation
  & Un échange écrit entre un client et le service client. \\[4pt]
Proposition de prix
  & Un prix négocié, proposé dans une conversation. Elle a une date
    d'expiration. \\[4pt]
Profil
  & Un paquet de droits que l'on attribue à un responsable. \\[4pt]
Service
  & Le regroupement métier d'un responsable, avec son chef. \\[4pt]
Réclamation
  & Un problème signalé par un client après la vente. \\[4pt]
Retour
  & Une demande de renvoi de marchandise. \\
\bottomrule
\end{longtable}

\begin{attention}
\textbf{Ne confondez pas réclamation et retour.} Une réclamation est un
\emph{problème} qu'on signale. Un retour est une \emph{marchandise} qui revient.
Les deux écrans sont différents et les deux circuits aussi.
\end{attention}

% =============================================================================
\chapter{Formation du Client}
\label{ch:client}
% =============================================================================

\begin{roleBox}
\textbf{Qui vous êtes.} Vous achetez sur GARAH. Vous n'avez aucun droit
particulier à demander : votre accès repose sur le fait que \emph{ce sont vos
données}. Vous voyez vos commandes, vos discussions, vos retours — et ceux de
personne d'autre.

\medskip
\textbf{Ce que vous pouvez faire}
\begin{itemize}[leftmargin=*,itemsep=2pt,topsep=3pt]
  \item parcourir le catalogue \textbf{sans compte} et remplir votre panier ;
  \item créer un compte et confirmer votre adresse e-mail ;
  \item commander, choisir votre point de récupération, payer par Mobile Money ;
  \item suivre un colis, \textbf{même sans compte}, avec son numéro de suivi ;
  \item poser une question sur un article et négocier un prix ;
  \item demander un retour ou déposer une réclamation ;
  \item tenir une liste d'articles mis de côté.
\end{itemize}

\medskip
\textbf{Ce que vous ne pouvez pas faire}
\begin{itemize}[leftmargin=*,itemsep=2pt,topsep=3pt]
  \item voir le back-office, ni les commandes des autres clients ;
  \item modifier un prix, un stock ou un statut de commande ;
  \item annuler une commande déjà payée vous-même — cela passe par un
        administrateur ;
  \item prendre en charge votre propre conversation.
\end{itemize}

\medskip
\textbf{Avec qui vous interagissez.} Le \textbf{responsable} du service client
répond à vos questions. Le \textbf{point de récupération} vous remet la
marchandise contre votre code.
\end{roleBox}

% -----------------------------------------------------------------------------
\section{Module 1 — Créer son compte et se connecter}

\subsection{Tutoriel : créer un compte}

\begin{objectif}
Vous allez créer votre compte client et confirmer votre adresse e-mail, sans
quoi certaines actions vous seront refusées plus tard.
\end{objectif}

\begin{avant}
Il vous faut une adresse e-mail à laquelle vous avez réellement accès, et un
mot de passe d'au moins huit caractères.
\end{avant}

\etape{Ouvrir la page d'inscription}

Sur la boutique, cliquez sur \textbf{« S'inscrire »} en haut à droite. Dans
l'application Android, ouvrez l'onglet \textbf{Compte}, puis
\textbf{« Créer un compte »}.

\begin{illustration}
\begin{center}
\begin{tikzpicture}[x=1mm,y=1mm]
  % La barre du site
  \draw[fill=white,draw=garahBordure,line width=0.8pt,rounded corners=2pt]
    (0,0) rectangle (150,14);
  \node[font=\scriptsize\bfseries,text=garahEncre] at (14,7) {GARAH};
  \node[font=\tiny,text=garahSlogan] at (16,3.4) {AU-DELÀ DES FRONTIÈRES};
  \node[font=\tiny,text=garahAttenue] at (44,7) {Accueil};
  \node[font=\tiny,text=garahAttenue] at (58,7) {Catalogue};
  \node[font=\tiny,text=garahAttenue] at (76,7) {Suivre un colis};
  \draw[fill=garahPrimaire,draw=none,rounded corners=2pt] (104,3) rectangle (124,11);
  \node[font=\tiny,text=white] at (114,7) {Se connecter};
  \draw[fill=white,draw=garahBordure,rounded corners=2pt] (126,3) rectangle (146,11);
  \node[font=\tiny,text=garahEncre] at (136,7) {S'inscrire};
  \node at (150,7) {\num{1}};
\end{tikzpicture}
\end{center}
\end{illustration}
\legendeIllus{2.1}{La barre du site. \num{1} le bouton « S'inscrire ».}

\etape{Remplir le formulaire}

Renseignez votre nom, votre prénom, votre adresse e-mail et votre mot de passe.

\begin{astuce}
Le \textbf{numéro de téléphone est facultatif} à l'inscription — mais
renseignez-le tout de suite. Il sera pré-rempli au moment de payer, et vous
n'aurez pas à le retaper. Le format est souple : \texttt{+237 6 99 00 00 00},
\texttt{0699000000} ou \texttt{(237) 699-000-000} sont tous acceptés.
\end{astuce}

\etape{Confirmer votre adresse e-mail}

Un message vous est envoyé. Ouvrez-le et cliquez sur le lien qu'il contient.

\begin{attention}
Tant que votre adresse n'est pas confirmée, la mention
\textbf{« NON VÉRIFIÉE »} apparaît à côté de votre compte. Ne remettez pas
cette étape à plus tard : certaines opérations la contrôlent, et vous
découvririez le refus au pire moment.
\end{attention}

\begin{resultat}
Vous êtes connecté. Votre nom apparaît en haut à droite de la boutique, ou dans
l'onglet \textbf{Compte} de l'application.
\end{resultat}

\subsection{Que faire si ça ne fonctionne pas ?}

\begin{longtable}{@{}p{3.7cm} p{3.7cm} p{6.8cm}@{}}
\toprule
\bfseries Problème & \bfseries Cause possible & \bfseries Solution \\
\midrule
\endhead
« Ce numéro de téléphone n'est pas valide »
  & Le champ contient des lettres, ou plus de 30 caractères
  & N'utilisez que des chiffres, espaces, points, tirets, parenthèses et le
    signe~+ \\[4pt]
Le message de confirmation n'arrive pas
  & L'adresse est mal saisie, ou le message est dans les indésirables
  & Vérifiez le dossier « spam ». Si l'adresse est fausse, recréez un compte
    avec la bonne \\[4pt]
« Mauvais identifiants » à la connexion
  & Mot de passe ou adresse incorrects
  & Le message ne dit pas lequel des deux est faux, et c'est volontaire : le
    préciser aiderait quelqu'un qui cherche à deviner \\[4pt]
Vous êtes déconnecté sans avoir rien fait
  & Votre session a expiré
  & Reconnectez-vous. Sur l'application, la session est conservée plus
    longtemps que sur le site \\
\bottomrule
\end{longtable}

% -----------------------------------------------------------------------------
\section{Module 2 — Trouver un article}

\subsection{Tutoriel : parcourir le catalogue}

\begin{objectif}
Vous allez apprendre les trois façons de trouver un article : les catégories,
la recherche, et le catalogue complet.
\end{objectif}

\etape{Utiliser les catégories}

Sur l'accueil, les catégories sont présentées sous forme de puces. Touchez-en
une : le catalogue s'ouvre \textbf{déjà filtré}.

\begin{illustration}
\begin{center}
\begin{tikzpicture}[x=1mm,y=1mm]
  \node[font=\scriptsize\bfseries,text=garahAttenue,anchor=west] at (0,26) {CATÉGORIES};
  \foreach \x/\t [count=\i] in {0/Chaussure, 26/Vetements, 55/Cheveux et beaute,
                                92/Informatique, 120/Electromenager} {
    \draw[draw=garahBordure,fill=white,rounded corners=4pt,line width=0.7pt]
      (\x,14) rectangle (\x+24,22);
    \node[font=\tiny,text=garahEncre] at (\x+12,18) {\t};
  }
  \node at (38,10) {\num{1}};
  \draw[-{Stealth[length=2mm]},garahPrimaire,line width=0.8pt] (38,12) -- (36,14);
\end{tikzpicture}
\end{center}
\end{illustration}
\legendeIllus{2.2}{Les puces de catégories sur l'accueil. \num{1} un clic filtre
directement le catalogue.}

\etape{Reconnaître le filtre actif}

Une fois dans le catalogue, le filtre est \textbf{visible} : une puce porte le
nom de la catégorie, avec une croix pour l'enlever.

\begin{attention}
Si vous ne voyez que quelques articles, \textbf{regardez d'abord s'il y a une
puce de filtre}. Un catalogue filtré ressemble à un catalogue vide.
\end{attention}

\etape{Chercher par mot}

La barre de recherche porte sur \textbf{le nom de l'article, le vendeur et la
catégorie}. Elle ne cherche pas dans les références ni dans les numéros de
commande.

\begin{astuce}
La recherche attend que vous ayez fini de taper avant d'interroger le serveur.
Ce n'est pas de la lenteur : c'est ce qui évite d'envoyer neuf requêtes pour le
mot « chaussure », sur un forfait qui se compte.
\end{astuce}

\subsection{Tutoriel : lire une fiche article}

\begin{objectif}
Comprendre les prix par palier, les déclinaisons, et le sous-total.
\end{objectif}

\etape{Comprendre les paliers de prix}

Le prix baisse par tranches de quantité. La grille est affichée
\textbf{avant} le bouton d'achat, et le palier qui s'applique est mis en avant.

\begin{illustration}
\begin{center}
\begin{tikzpicture}[x=1mm,y=1mm]
  \node[font=\scriptsize\bfseries,text=garahAttenue,anchor=west] at (0,32) {PRIX PAR QUANTITÉ};
  % Palier actif
  \draw[draw=garahPrimaire,fill=garahPrimaire!8,rounded corners=3pt,line width=1.1pt]
    (0,10) rectangle (36,28);
  \node[font=\tiny,text=garahAttenue] at (18,25) {1 -- 6};
  \node[font=\scriptsize\bfseries,text=garahEncre] at (18,19) {2 000 FCFA};
  \node[font=\tiny,text=garahAttenue] at (18,14) {la pièce};
  % Autres paliers
  \draw[draw=garahBordure,fill=white,rounded corners=3pt,line width=0.8pt]
    (40,10) rectangle (76,28);
  \node[font=\tiny,text=garahAttenue] at (58,25) {7 -- 10};
  \node[font=\scriptsize\bfseries,text=garahEncre] at (58,19) {1 500 FCFA};
  \node[font=\tiny,text=garahAttenue] at (58,14) {la pièce};
  \draw[draw=garahBordure,fill=white,rounded corners=3pt,line width=0.8pt]
    (80,10) rectangle (116,28);
  \node[font=\tiny,text=garahAttenue] at (98,25) {11 et +};
  \node[font=\scriptsize\bfseries,text=garahEncre] at (98,19) {1 000 FCFA};
  \node[font=\tiny,text=garahAttenue] at (98,14) {la pièce};

  \node[font=\tiny,text=garahAttenue,anchor=west] at (0,5)
    {Encore \textbf{6} et le prix passe à \textbf{1 500 FCFA}.};
  \node at (-6,19) {\num{1}};
  \node at (122,19) {\num{2}};
\end{tikzpicture}
\end{center}
\end{illustration}
\legendeIllus{2.3}{La grille des paliers. \num{1} le palier qui s'applique à la
quantité choisie, encadré. \num{2} les paliers suivants, atteignables en
commandant davantage. La phrase du bas indique ce qu'il manque pour descendre
d'un palier.}

\etape{Choisir la déclinaison}

Si l'article existe en plusieurs versions, choisissez la vôtre. Les
déclinaisons \textbf{barrées sont épuisées} — elles restent visibles pour que
vous sachiez qu'elles existent.

\etape{Vérifier le sous-total}

Le sous-total s'affiche sous la quantité.

\begin{attention}
Le sous-total affiché sur la fiche \textbf{n'est pas ce que vous paierez} : il
ne contient pas les frais d'acheminement, qui dépendent du point de
récupération — et vous ne l'avez pas encore choisi.

Le montant définitif est \textbf{figé au moment de la commande}.
\end{attention}

\begin{astuce}
Si le bouton « + » de la quantité cesse de répondre, lisez le message juste
au-dessous. Deux raisons possibles, et elles n'appellent pas le même geste :
\begin{itemize}[leftmargin=*,itemsep=1pt,topsep=2pt]
  \item \textbf{le stock} est atteint — il n'y en a plus ;
  \item \textbf{le prix n'est pas fixé} au-delà — écrivez au marchand avec le
        bouton « Contacter ».
\end{itemize}
\end{astuce}

\begin{resultat}
Vous voyez le bon prix unitaire, la bonne quantité, et un sous-total cohérent.
Si l'écran affiche « Prix non fixé pour cette quantité », le bouton d'ajout est
désactivé : c'est normal, aucun prix n'existe pour cette quantité-là.
\end{resultat}

% -----------------------------------------------------------------------------
\section{Module 3 — Commander et payer}

\subsection{Tutoriel : passer une commande}

\begin{objectif}
De l'ajout au panier jusqu'au paiement confirmé.
\end{objectif}

\begin{avant}
Votre panier contient au moins un article. Vous connaissez le numéro Mobile
Money qui recevra la demande de validation.
\end{avant}

\etape{Ajouter au panier}

Sur la fiche, choisissez la quantité puis touchez \textbf{« Ajouter au
panier »}.

\begin{itemize}[leftmargin=*,itemsep=2pt]
  \item Sur \textbf{l'application}, un bandeau confirme l'ajout et propose
        « Voir le panier ».
  \item Sur \textbf{la boutique}, le bouton devient un lien « Voir le panier ».
\end{itemize}

\begin{astuce}
Vous n'avez \textbf{pas besoin de compte} pour remplir votre panier. La
connexion est demandée seulement au moment de commander.
\end{astuce}

\etape{Ouvrir le panier et vérifier}

Contrôlez les quantités. Le sous-total est affiché, sans les frais
d'acheminement.

\etape{Passer commande}

Touchez \textbf{« Passer commande »}. Si vous n'êtes pas connecté, l'écran de
connexion s'ouvre d'abord, et votre panier est repris automatiquement après.

\etape{Choisir le point de récupération}

C'est ici que les \textbf{frais d'acheminement} apparaissent : ils dépendent du
point choisi. Le moins cher est proposé par défaut.

\begin{illustration}
\begin{center}
\begin{tikzpicture}[x=1mm,y=1mm]
  \draw[draw=garahBordure,fill=white,rounded corners=3pt,line width=0.8pt]
    (0,0) rectangle (130,44);
  \node[font=\scriptsize\bfseries,text=garahAttenue,anchor=west] at (5,39)
    {POINT DE RÉCUPÉRATION};
  \draw[draw=garahPrimaire,fill=garahPrimaire!7,rounded corners=3pt,line width=1pt]
    (5,26) rectangle (125,35);
  \filldraw[garahPrimaire] (10,30.5) circle (1.6);
  \node[font=\tiny,text=garahEncre,anchor=west] at (14,30.5) {Point A — frais les moins élevés};
  \draw[draw=garahBordure,fill=white,rounded corners=3pt,line width=0.7pt]
    (5,14) rectangle (125,23);
  \draw[garahAttenue,line width=0.7pt] (10,18.5) circle (1.6);
  \node[font=\tiny,text=garahEncre,anchor=west] at (14,18.5) {Point B};
  \node[font=\tiny,text=garahAttenue,anchor=west] at (5,8)
    {Articles \quad Frais d'acheminement \quad \textbf{Total à payer}};
  \node at (132,30.5) {\num{1}};
  \node at (132,6) {\num{2}};
\end{tikzpicture}
\end{center}
\end{illustration}
\legendeIllus{2.4}{L'écran de commande. \num{1} le point de récupération choisi.
\num{2} le total, qui ne se complète qu'une fois le point sélectionné.}

\etape{Payer par Mobile Money}

Choisissez \textbf{MTN Mobile Money} ou \textbf{Orange Money}, vérifiez le
numéro, puis lancez le paiement.

\begin{attention}
Le numéro est \textbf{pré-rempli depuis votre profil, mais modifiable}. C'est
volontaire : on paie souvent avec un autre numéro que celui du compte — celui
d'un proche, ou son second opérateur.

Vérifiez-le avant de valider : c'est \emph{ce} numéro qui recevra la demande.
\end{attention}

\etape{Valider sur votre téléphone}

Une demande de validation arrive sur le numéro indiqué. Acceptez-la.

\begin{attention}
L'écran \textbf{n'affiche pas « payé » tant que l'opérateur n'a pas confirmé}.
Ce n'est pas de la lenteur : afficher « payé » trop tôt ferait repartir
quelqu'un persuadé d'avoir réglé, et le litige coûterait plus cher que
l'attente.

Quand vous avez validé sur votre téléphone, appuyez sur le bouton qui redemande
l'état. \textbf{L'écran ne se rafraîchit pas tout seul}, et c'est délibéré :
interroger en boucle viderait votre batterie et votre forfait.
\end{attention}

\begin{resultat}
Le message \textbf{« Paiement confirmé »} s'affiche. Votre commande est
enregistrée, votre panier est vidé, et vous serez prévenu lorsque la
marchandise sera disponible au point de récupération.
\end{resultat}

\subsection{Comprendre l'état de votre commande}

Une commande passe par des états successifs. Voici lesquels, et ce qu'ils
veulent dire pour vous :

\begin{illustration}
\begin{center}
\begin{tikzpicture}[
  node distance=4mm,
  et/.style={rectangle,rounded corners=3pt,draw=garahMarque,fill=garahMarque!8,
             line width=0.8pt,minimum height=9mm,text width=27mm,align=center,
             font=\scriptsize\bfseries},
  ko/.style={rectangle,rounded corners=3pt,draw=garahDanger,fill=garahDanger!8,
             line width=0.8pt,minimum height=9mm,text width=27mm,align=center,
             font=\scriptsize\bfseries},
  f/.style={-{Stealth[length=2.6mm]},garahAttenue,line width=0.9pt}
]
\node[et] (a) {En attente\\de paiement};
\node[et,right=of a] (b) {Payée};
\node[et,right=of b] (c) {En préparation};
\node[et,right=of c] (d) {Prête};
\node[et,below=8mm of d] (e) {Expédiée};
\node[et,left=of e] (f) {Disponible};
\node[et,left=of f] (g) {Retirée};
\node[ko,below=8mm of a] (x) {Annulée};

\draw[f] (a) -- (b); \draw[f] (b) -- (c); \draw[f] (c) -- (d);
\draw[f] (d) -- (e); \draw[f] (e) -- (f); \draw[f] (f) -- (g);
\draw[f] (a) -- (x);
\draw[f] (b) -- ($(b)+(0,-8mm)$) -- (x);
\end{tikzpicture}
\end{center}
\end{illustration}
\legendeIllus{2.5}{Le cycle de vie d'une commande. Une commande non payée peut
être annulée par le client ou par le délai ; une commande payée ne peut l'être
que par un administrateur, avec remboursement.}

\begin{center}
\begin{tabular}{@{}>{\bfseries}p{3.6cm} p{11.2cm}@{}}
\toprule
État & Ce que vous devez faire \\
\midrule
En attente de paiement & Payer. Sans quoi la commande sera annulée. \\
Payée & Rien. L'équipe prend le relais. \\
En préparation & Rien. Votre marchandise est rassemblée. \\
Prête & Rien. Elle attend son départ. \\
Expédiée & Rien. Elle voyage. \\
Disponible & \textbf{Aller la retirer}, avec votre code. \\
Retirée & C'est terminé. \\
\bottomrule
\end{tabular}
\end{center}

% -----------------------------------------------------------------------------
\section{Module 4 — Suivre et retirer}

\subsection{Tutoriel : suivre un colis}

\begin{objectif}
Savoir où en est votre marchandise, avec ou sans compte.
\end{objectif}

\etape{Ouvrir le suivi}

Menu \textbf{« Suivre un colis »} sur la boutique, ou l'écran de suivi dans
l'application.

\begin{astuce}
Le suivi fonctionne \textbf{sans compte}. Il suffit du numéro de suivi — vous
pouvez donc le transmettre à la personne qui viendra retirer à votre place.
\end{astuce}

\etape{Saisir le numéro de suivi}

Entrez le numéro. L'état du colis s'affiche.

\begin{illustration}
\begin{center}
\begin{tikzpicture}[
  node distance=5mm,
  et/.style={rectangle,rounded corners=3pt,draw=garahMarque,fill=garahMarque!8,
             line width=0.8pt,minimum height=8mm,text width=24mm,align=center,
             font=\scriptsize\bfseries},
  bl/.style={rectangle,rounded corners=3pt,draw=garahAlerte,fill=garahAlerte!12,
             line width=0.8pt,minimum height=8mm,text width=24mm,align=center,
             font=\scriptsize\bfseries},
  f/.style={-{Stealth[length=2.6mm]},garahAttenue,line width=0.9pt}
]
\node[et] (a) {Créé};
\node[et,right=of a] (b) {En transit};
\node[et,right=of b] (c) {Disponible};
\node[et,right=of c] (d) {Remis};
\node[bl,below=7mm of b] (e) {Bloqué};
\draw[f] (a) -- (b); \draw[f] (b) -- (c); \draw[f] (c) -- (d);
\draw[f] (b) -- (e);
\end{tikzpicture}
\end{center}
\end{illustration}
\legendeIllus{2.6}{Les états d'un colis. « Bloqué » signale un incident en
cours de route : contactez le service client.}

\subsection{Tutoriel : retirer sa marchandise}

\begin{objectif}
Récupérer votre commande au point de récupération.
\end{objectif}

\begin{avant}
Votre commande doit être à l'état \textbf{Disponible}. Vous devez avoir votre
\textbf{code de retrait}.
\end{avant}

\etape{Attendre l'avis de disponibilité}

Vous êtes prévenu lorsque la marchandise arrive au point.

\etape{Se présenter avec le code}

Donnez votre code de retrait à l'agent du point.

\begin{attention}
\textbf{Le code de retrait est le seul moyen de prouver la remise.} Sans lui,
votre parole s'oppose à celle du point de récupération et personne ne peut
trancher.

Ne le communiquez qu'à la personne que vous chargez de retirer à votre place.
\end{attention}

\begin{resultat}
La marchandise vous est remise et la commande passe à l'état
\textbf{Retirée}.
\end{resultat}

% -----------------------------------------------------------------------------
\section{Module 5 — Dialoguer avec GARAH}

\subsection{Tutoriel : poser une question sur un article}

\begin{objectif}
Ouvrir une conversation depuis une fiche article, et suivre la réponse.
\end{objectif}

\etape{Ouvrir la conversation}

Sur la fiche de l'article, touchez \textbf{« Contacter »}. Une conversation est
créée, avec l'article en objet.

\etape{Attendre qu'un conseiller la prenne}

L'écran indique \textbf{« En attente d'un conseiller »} tant que personne ne
s'en est chargé, puis \textbf{« Un conseiller vous suit »}.

\begin{attention}
Tant que la conversation est \textbf{en attente}, personne n'y est encore
affecté. Vos messages sont bien enregistrés, mais la réponse en direct
n'arrivera qu'une fois la conversation prise en charge.
\end{attention}

\etape{Échanger}

Écrivez dans le champ du bas. \textbf{Les messages arrivent en direct}, dans
les deux sens : vous n'avez pas à recharger la page.

\begin{illustration}
\begin{center}
\begin{tikzpicture}[x=1mm,y=1mm]
  % Bulle client (droite)
  \draw[draw=garahPrimaire!40,fill=garahPrimaire!12,rounded corners=3pt]
    (52,26) rectangle (120,36);
  \node[font=\tiny,text=garahEncre,anchor=west] at (55,32) {Bonjour, je voudrais en savoir plus.};
  \node[font=\tiny,text=garahAttenue,anchor=west] at (55,28) {10 sept., 00:39};
  % Bulle conseiller (gauche)
  \draw[draw=garahBordure,fill=garahFond,rounded corners=3pt]
    (0,12) rectangle (68,22);
  \node[font=\tiny,text=garahEncre,anchor=west] at (3,18) {Bonjour, cette chaussure coûte 2 000 F.};
  \node[font=\tiny,text=garahAttenue,anchor=west] at (3,14) {10 sept., 00:42};
  % Champ
  \draw[draw=garahBordure,fill=white,rounded corners=2pt] (0,0) rectangle (98,8);
  \node[font=\tiny,text=garahAttenue,anchor=west] at (3,4) {Votre message…};
  \draw[fill=garahPrimaire,draw=none,rounded corners=2pt] (102,0) rectangle (120,8);
  \node[font=\tiny,text=white] at (111,4) {Envoyer};
  \node at (126,31) {\num{1}};
  \node at (74,17) {\num{2}};
  \node at (126,4) {\num{3}};
\end{tikzpicture}
\end{center}
\end{illustration}
\legendeIllus{2.7}{Une conversation. \num{1} vos messages, à droite. \num{2} les
réponses du conseiller, à gauche. \num{3} le champ de saisie.}

\subsection{Comprendre une proposition de prix}

Un conseiller peut vous faire une \textbf{proposition de prix} dans la
conversation.

\begin{attention}
Une proposition \textbf{expire}. Passé sa date, elle devient « Expirée » et ne
peut plus être acceptée — il faut en redemander une. Ne la laissez pas dormir.
\end{attention}

\begin{illustration}
\begin{center}
\begin{tikzpicture}[
  node distance=6mm,
  et/.style={rectangle,rounded corners=3pt,line width=0.8pt,
             minimum height=8mm,text width=24mm,align=center,
             font=\scriptsize\bfseries},
  f/.style={-{Stealth[length=2.6mm]},garahAttenue,line width=0.9pt}
]
\node[et,draw=garahMarque,fill=garahMarque!8] (a) {Proposée};
\node[et,draw=garahSucces,fill=garahSucces!10,right=14mm of a] (b) {Acceptée};
\node[et,draw=garahDanger,fill=garahDanger!8,below=5mm of b] (c) {Refusée};
\node[et,draw=garahAlerte,fill=garahAlerte!12,below=5mm of c] (d) {Expirée};
\draw[f] (a) -- (b);
\draw[f] (a) -- (c);
\draw[f] (a) -- (d);
\end{tikzpicture}
\end{center}
\end{illustration}
\legendeIllus{2.8}{Les issues d'une proposition de prix.}

% -----------------------------------------------------------------------------
\section{Module 6 — Après la vente}

\subsection{Demander un retour}

Un retour concerne \textbf{une marchandise qui revient}. Il suit ce circuit :

\begin{illustration}
\begin{center}
\begin{tikzpicture}[
  node distance=5mm,
  et/.style={rectangle,rounded corners=3pt,line width=0.8pt,
             minimum height=8mm,text width=22mm,align=center,
             font=\scriptsize\bfseries},
  f/.style={-{Stealth[length=2.6mm]},garahAttenue,line width=0.9pt}
]
\node[et,draw=garahMarque,fill=garahMarque!8] (a) {Demandé};
\node[et,draw=garahSucces,fill=garahSucces!10,right=of a] (b) {Accepté};
\node[et,draw=garahMarque,fill=garahMarque!8,right=of b] (c) {Réceptionné};
\node[et,draw=garahSucces,fill=garahSucces!10,right=of c] (d) {Validé};
\node[et,draw=garahDanger,fill=garahDanger!8,below=6mm of b] (e) {Refusé};
\draw[f] (a) -- (b); \draw[f] (b) -- (c); \draw[f] (c) -- (d);
\draw[f] (a) -- (e);
\end{tikzpicture}
\end{center}
\end{illustration}
\legendeIllus{2.9}{Le circuit d'un retour. Chaque étape est décidée par
l'équipe ; vous suivez l'avancement depuis « Mes retours ».}

\subsection{Déposer une réclamation}

Une réclamation concerne \textbf{un problème}, pas une marchandise. Elle passe
de \textbf{Ouverte} à \textbf{En cours} lorsqu'un responsable s'en saisit, puis
à \textbf{Résolue}.

\begin{retenir}
\textbf{Retour} = quelque chose revient. \textbf{Réclamation} = quelque chose ne
va pas. Deux écrans, deux circuits, ne les confondez pas.
\end{retenir}

% -----------------------------------------------------------------------------
\section{Scénario complet : de la question à la marchandise}

\begin{tcolorbox}[colback=garahFond,colframe=garahEncre,arc=4pt,boxrule=1pt]
\textbf{Situation.} Vous cherchez des chaussures pour votre commerce. Vous en
voulez une dizaine, et vous voulez savoir si le prix peut baisser.
\end{tcolorbox}

\begin{enumerate}[leftmargin=*,itemsep=4pt]
  \item Vous ouvrez la boutique et touchez la catégorie \textbf{Chaussure}.
  \item Vous ouvrez la fiche et lisez la grille : à partir de 7 pièces, le prix
        passe de 2 000 à 1 500 FCFA.
  \item Vous touchez \textbf{« Contacter »} et demandez un prix pour 20 pièces.
  \item Le conseiller prend la conversation — la mention passe à
        \textbf{« Un conseiller vous suit »}.
  \item Il vous fait une proposition de prix. Vous l'acceptez \emph{avant son
        expiration}.
  \item Vous ajoutez la quantité au panier, puis \textbf{« Passer commande »}.
  \item Vous choisissez votre point de récupération : les frais s'affichent, le
        total se complète.
  \item Vous payez par Mobile Money, validez sur votre téléphone, et appuyez sur
        le bouton qui redemande l'état.
  \item Vous suivez la commande jusqu'à \textbf{Disponible}.
  \item Vous vous présentez au point avec votre \textbf{code de retrait}.
\end{enumerate}

\begin{exercice}
\textbf{Objectif.} Passer une commande complète sans aide.

\medskip
\textbf{Actions attendues}
\begin{enumerate}[leftmargin=*,itemsep=1pt,topsep=2pt]
  \item Créer un compte et confirmer l'adresse e-mail.
  \item Trouver un article \emph{par la catégorie}, pas par la recherche.
  \item Choisir une quantité qui déclenche le \emph{deuxième} palier de prix.
  \item Ajouter au panier et vérifier le sous-total.
  \item Passer commande, choisir un point de récupération.
  \item Aller jusqu'à l'écran de paiement \emph{sans le valider}.
\end{enumerate}

\medskip
\textbf{Résultat attendu.} L'écran de paiement affiche le numéro pré-rempli et
un total qui inclut les frais d'acheminement.

\medskip
\textbf{Critères de réussite}
\begin{itemize}[leftmargin=*,itemsep=1pt,topsep=2pt]
  \item le prix unitaire correspond bien au deuxième palier ;
  \item le total est supérieur au sous-total du panier (les frais s'ajoutent) ;
  \item vous savez dire pourquoi.
\end{itemize}
\end{exercice}

\begin{checklist}
\case Je sais créer un compte et confirmer mon adresse\\
\case Je sais me connecter et me déconnecter\\
\case Je sais trouver un article par catégorie et par recherche\\
\case Je sais reconnaître un filtre actif et l'enlever\\
\case Je sais lire une grille de paliers et dire quel prix s'applique\\
\case Je sais distinguer le sous-total du total à payer\\
\case Je sais ajouter au panier et vérifier mon panier\\
\case Je sais choisir un point de récupération\\
\case Je sais payer par Mobile Money et vérifier la confirmation\\
\case Je sais suivre un colis, avec et sans compte\\
\case Je sais retirer ma marchandise avec mon code\\
\case Je sais ouvrir une conversation et répondre\\
\case Je sais ce qu'est une proposition de prix, et qu'elle expire\\
\case Je sais distinguer un retour d'une réclamation\\
\case Je sais quoi faire si le paiement ne se confirme pas
\end{checklist}

% =============================================================================
\chapter{Formation du Responsable}
\label{ch:responsable}
% =============================================================================

\begin{roleBox}
\textbf{Qui vous êtes.} Vous faites le travail quotidien de GARAH depuis le
back-office : répondre aux clients, préparer les commandes, suivre les
expéditions, confirmer les retraits, traiter l'après-vente.

\medskip
\textbf{Ce que vous pouvez faire.} \emph{Cela dépend de vos profils.} Un
responsable n'a aucun droit en propre : un administrateur lui attribue un ou
plusieurs \textbf{profils}, et chaque profil est un paquet de droits.

\begin{attention}
\textbf{Si un écran n'apparaît pas dans votre menu, ce n'est pas une panne.}
C'est que vos profils ne vous en donnent pas l'accès. Le menu ne montre que ce
que vous pouvez réellement ouvrir — un back-office qui afficherait vingt écrans
dont quinze répondent « accès refusé » serait illisible.

Pour obtenir un accès, demandez-le à un administrateur. N'essayez pas de taper
l'adresse à la main : le serveur refusera de toute façon.
\end{attention}

\medskip
\textbf{Ce que vous ne pouvez pas faire}
\begin{itemize}[leftmargin=*,itemsep=2pt,topsep=3pt]
  \item créer ou modifier des comptes de l'équipe, ni attribuer des profils ;
  \item annuler une commande déjà payée (réservé à l'administrateur) ;
  \item consulter le journal des actions (réservé au super administrateur) ;
  \item voir les dossiers de toute l'équipe si votre filtre est sur
        « Mes dossiers ».
\end{itemize}

\medskip
\textbf{Avec qui vous interagissez.} Les \textbf{clients}, par les
conversations et l'après-vente. Vos \textbf{collègues}, par la messagerie
interne. Votre \textbf{administrateur}, pour vos droits.
\end{roleBox}

% -----------------------------------------------------------------------------
\section{Module 1 — Se connecter et s'orienter}

\subsection{Tutoriel : votre première connexion}

\begin{objectif}
Vous connecter au back-office et comprendre l'organisation du menu.
\end{objectif}

\begin{avant}
Un administrateur a créé votre compte et vous a attribué au moins un profil.
Vous avez confirmé votre adresse e-mail.
\end{avant}

\etape{Se connecter}

Ouvrez l'adresse du back-office et saisissez votre adresse e-mail et votre mot
de passe.

\etape{Reconnaître le menu}

Le menu de gauche est rangé dans \textbf{l'ordre de la vie d'une vente}, en
sept groupes. Ce n'est pas une liste alphabétique : on cherche d'abord la bonne
\emph{question}, puis l'écran dedans.

\begin{illustration}
\begin{center}
\begin{tikzpicture}[x=1mm,y=1mm]
  \draw[fill=garahSurfaceNuit,draw=garahAccent!35,line width=0.8pt,rounded corners=3pt]
    (0,0) rectangle (58,132);

  % Marque
  \node[font=\small\bfseries,text=white,anchor=west] at (5,125) {GARAH};
  \node[font=\tiny,text=garahSucces,anchor=west] at (5,120.5) {AU-DELÀ DES FRONTIÈRES};
  \node[font=\tiny,text=garahAttenue,anchor=west] at (5,116.5) {ADMINISTRATION};

  \def\grp#1#2{\node[font=\tiny\bfseries,text=garahAttenue,anchor=west] at (5,#1) {#2};}
  \def\ent#1#2{\node[font=\scriptsize,text=white!85,anchor=west] at (9,#1) {#2};
               \draw[garahMarque!60,fill=garahMarque!30] (6,#1) circle (0.9);}

  \grp{108}{PILOTAGE}
  \ent{102}{Tableau de bord} \ent{96}{Statistiques}

  \grp{88}{LE CATALOGUE}
  \ent{82}{Produits} \ent{76}{Catégories} \ent{70}{Dimensions}
  \ent{64}{Stock} \ent{58}{Marchands}

  \grp{50}{LES VENTES}
  \ent{44}{Commandes} \ent{38}{Paiements} \ent{32}{Marchands à payer}

  \grp{24}{L'ACHEMINEMENT}
  \ent{18}{Expéditions} \ent{12}{Retraits} \ent{6}{Lieux, Itinéraires}

  \node at (64,108) {\num{1}};
  \node at (64,88)  {\num{2}};
  \node at (64,50)  {\num{3}};
  \node at (64,24)  {\num{4}};

  % Les trois groupes restants, en encart
  \draw[draw=garahBordure,fill=garahFond,rounded corners=3pt,line width=0.7pt]
    (72,60) rectangle (140,110);
  \node[font=\tiny\bfseries,text=garahAttenue,anchor=west] at (76,104) {APRÈS-VENTE};
  \node[font=\scriptsize,anchor=west] at (78,99) {Réclamations};
  \node[font=\scriptsize,anchor=west] at (78,94) {Retours};
  \node[font=\scriptsize,anchor=west] at (78,89) {Service client};
  \node[font=\tiny\bfseries,text=garahAttenue,anchor=west] at (76,82) {LES GENS};
  \node[font=\scriptsize,anchor=west] at (78,77) {Clients, Équipe, Mon service};
  \node[font=\tiny\bfseries,text=garahAttenue,anchor=west] at (76,70) {LA TRACE};
  \node[font=\scriptsize,anchor=west] at (78,65) {Journal des actions};
  \node at (146,85) {\num{5}};
\end{tikzpicture}
\end{center}
\end{illustration}
\legendeIllus{3.1}{Le menu du back-office, rangé par question.
\num{1} où en est-on ? \num{2} ce qu'on vend. \num{3} ce qui se vend.
\num{4} ce qui part. \num{5} l'après-vente, les gens, et la trace.
Les entrées absentes de \emph{votre} menu correspondent à des droits que vous
n'avez pas.}

\captureAInserer[10em]{Le menu réel de votre installation — il varie selon vos
profils.}

\begin{astuce}
Un chiffre à côté d'une entrée du menu indique un \textbf{nombre en attente}.
Il n'apparaît que s'il est supérieur à zéro : un « 0 » permanent cesse d'être
lu au bout de deux jours, et le jour où il devient « 1 », personne ne le voit.
\end{astuce}

\etape{Repérer votre compte}

En bas de la barre latérale : votre nom, votre rôle, et le bouton de
déconnexion. C'est aussi par là qu'on accède à son profil.

% -----------------------------------------------------------------------------
\section{Module 2 — Le service client}

C'est le cœur du travail quotidien. Une conversation passe par trois états :

\begin{illustration}
\begin{center}
\begin{tikzpicture}[
  node distance=16mm,
  et/.style={rectangle,rounded corners=3pt,line width=1pt,
             minimum height=11mm,text width=30mm,align=center,
             font=\small\bfseries},
  f/.style={-{Stealth[length=3mm]},garahAttenue,line width=1pt},
  lg/.style={font=\scriptsize\itshape,text=garahAttenue}
]
\node[et,draw=garahAlerte,fill=garahAlerte!12] (a) {En attente};
\node[et,draw=garahMarque,fill=garahMarque!10,right=of a] (b) {En cours};
\node[et,draw=garahAttenue,fill=garahFond,right=of b] (c) {Fermée};
\draw[f] (a) -- node[lg,above=1pt]{« Prendre »} (b);
\draw[f] (b) -- node[lg,above=1pt]{« Clore »} (c);
\node[lg,below=2mm of a,text width=32mm,align=center]
  {aucun responsable\\ affecté};
\node[lg,below=2mm of b,text width=32mm,align=center]
  {vous en êtes\\ responsable};
\node[lg,below=2mm of c,text width=32mm,align=center]
  {plus de réponse\\ possible};
\end{tikzpicture}
\end{center}
\end{illustration}
\legendeIllus{3.2}{Le cycle d'une conversation.}

\subsection{Tutoriel : prendre en charge une conversation}

\begin{objectif}
Repérer une demande en attente, la prendre, y répondre, puis la clore.
\end{objectif}

\begin{avant}
Vous devez disposer du droit de consulter les conversations, et de celui de les
prendre en charge.
\end{avant}

\etape{Ouvrir le service client}

Menu \textbf{Après-vente $\rightarrow$ Service client}.

\etape{Filtrer sur « En attente »}

Les onglets en haut filtrent la liste. \textbf{« En attente »} montre les
dossiers que personne n'a encore pris.

\begin{illustration}
\begin{center}
\begin{tikzpicture}[x=1mm,y=1mm]
  % Onglets
  \draw[fill=garahPrimaire,draw=none,rounded corners=8pt] (0,42) rectangle (18,50);
  \node[font=\tiny,text=white] at (9,46) {Toutes};
  \foreach \x/\t in {21/En attente, 50/En cours, 76/Fermée} {
    \draw[draw=garahBordure,fill=white,rounded corners=8pt] (\x,42) rectangle (\x+24,50);
    \node[font=\tiny,text=garahEncre] at (\x+12,46) {\t};
  }
  \node at (34,54) {\num{1}};
  \draw[-{Stealth[length=2mm]},garahPrimaire,line width=0.8pt] (34,52.6) -- (33,50.5);

  % Une carte de conversation
  \draw[draw=garahPrimaire!40,fill=white,rounded corners=3pt,line width=0.8pt]
    (0,6) rectangle (140,38);
  \draw[garahPrimaire,line width=1.6pt] (0.8,6) -- (0.8,38);
  \node[font=\scriptsize\bfseries,text=garahEncre,anchor=west] at (5,33)
    {Question : Chaussure Nike};
  \node[font=\tiny,text=garahEncre,anchor=west] at (5,28) {Dominique Bria};
  \node[font=\tiny,text=garahAttenue,anchor=west] at (28,28) {· CLI-000001};
  % Pastille non lus
  \draw[fill=garahPrimaire!14,draw=none,rounded corners=6pt] (5,16) rectangle (28,22);
  \node[font=\tiny,text=garahPrimaire] at (16.5,19) {2 non lu(s)};
  \node[font=\tiny,text=garahAttenue,anchor=west] at (32,19) {À l'instant};
  % Etiquette statut
  \draw[fill=garahAlerte!18,draw=none,rounded corners=6pt] (104,31) rectangle (136,37);
  \node[font=\tiny,text=garahAlerte!70!black] at (120,34) {EN ATTENTE};
  % Bouton Prendre
  \draw[fill=garahPrimaire,draw=none,rounded corners=2pt] (112,14) rectangle (136,22);
  \node[font=\tiny,text=white] at (124,18) {Prendre};
  \node at (146,18) {\num{3}};
  \node at (16.5,11) {\num{2}};
\end{tikzpicture}
\end{center}
\end{illustration}
\legendeIllus{3.3}{La file du service client. \num{1} l'onglet « En attente ».
\num{2} le compteur de messages non lus — il ne compte que les messages
\emph{du client}, jamais vos propres réponses. \num{3} le bouton « Prendre ».}

\etape{Prendre la conversation}

Cliquez sur \textbf{« Prendre »}. Le dossier vous est affecté et passe à
\textbf{En cours}.

\begin{attention}
La prise en charge est \textbf{atomique} : si un collègue clique une fraction de
seconde avant vous, c'est lui qui l'obtient et un message vous le dit. Ce n'est
pas une erreur — c'est ce qui empêche deux conseillers de répondre en même temps
au même client.
\end{attention}

\etape{Lire et répondre}

Ouvrez le dossier. Le fil s'affiche à droite, avec le champ de réponse.

\begin{attention}
\textbf{Les messages arrivent en direct, dans les deux sens.} Vous n'avez pas à
recharger. Si rien n'arrive alors que le client vous dit avoir écrit,
vérifiez que la conversation vous est bien \emph{affectée} : tant qu'elle est
« en attente », rien n'est poussé côté conseiller.
\end{attention}

\etape{Clore la conversation}

Quand la demande est traitée, cliquez sur \textbf{« Clore la conversation »}.

\begin{attention}
\textbf{Une conversation fermée n'accepte plus de message.} Le client devra en
ouvrir une nouvelle. Ne clôturez qu'une fois certain que le sujet est réglé.
\end{attention}

\begin{resultat}
Le dossier disparaît de « En attente », apparaît dans « En cours » puis dans
« Fermée ». Le compteur de non lus retombe à zéro.
\end{resultat}

\subsection{Que faire si ça ne fonctionne pas ?}

\begin{longtable}{@{}p{4cm} p{4cm} p{6.2cm}@{}}
\toprule
\bfseries Problème & \bfseries Cause possible & \bfseries Solution \\
\midrule
\endhead
« Prendre » ne fait rien
  & Un collègue a été plus rapide
  & Actualisez : le dossier est maintenant « En cours » à son nom \\[4pt]
Le message du client n'apparaît pas en direct
  & La conversation n'est pas encore affectée
  & Prenez-la d'abord. Sans responsable, rien n'est poussé \\[4pt]
Impossible de répondre
  & La conversation est fermée
  & Demandez au client d'en ouvrir une nouvelle \\[4pt]
« Service client » absent du menu
  & Vos profils ne donnent pas l'accès
  & Demandez le droit à un administrateur \\[4pt]
Le compteur de non lus semble faux
  & Il ne compte que les messages \emph{du client}
  & C'est le comportement attendu : vos réponses n'y figurent pas \\
\bottomrule
\end{longtable}

% -----------------------------------------------------------------------------
\section{Module 3 — Les commandes et l'acheminement}

\subsection{Tutoriel : suivre et faire avancer une commande}

\begin{objectif}
Retrouver une commande, comprendre son état, et la faire progresser.
\end{objectif}

\etape{Ouvrir les commandes}

Menu \textbf{Les ventes $\rightarrow$ Commandes}.

\etape{Comprendre l'état}

Reportez-vous à la figure~2.5 du chapitre client : le cycle est le même, vu de
l'autre côté. Ce qui vous concerne :

\begin{center}
\begin{tabular}{@{}>{\bfseries}p{3.8cm} p{11cm}@{}}
\toprule
État & Votre action \\
\midrule
En attente de paiement & Attendre. Ne préparez rien. \\
Payée & Lancer la préparation. \\
En préparation & Rassembler la marchandise. \\
Prête & Rattacher à une expédition. \\
Expédiée & Suivre le colis. \\
Disponible & Prévenir, puis confirmer le retrait. \\
Retirée & Terminé. \\
\bottomrule
\end{tabular}
\end{center}

\begin{attention}
\textbf{Ne préparez jamais une commande non payée.} L'état « En attente de
paiement » signifie que l'argent n'est pas arrivé. Une commande non payée peut
être annulée par le client ou par le délai.
\end{attention}

\subsection{Tutoriel : confirmer un retrait}

\begin{objectif}
Remettre la marchandise au client, contre son code.
\end{objectif}

\begin{avant}
La commande doit être \textbf{Disponible} au point. Le client doit vous
présenter son \textbf{code de retrait}.
\end{avant}

\etape{Ouvrir les retraits}

Menu \textbf{L'acheminement $\rightarrow$ Retraits}.

\etape{Vérifier le code}

Saisissez ou vérifiez le code présenté par le client.

\begin{attention}
\textbf{Le code de retrait est la seule preuve de la remise.} Sans lui, la
parole du client s'oppose à la vôtre et personne ne peut trancher.

Ne confirmez jamais un retrait « de confiance », sans code. C'est vous qui
serez interrogé en cas de litige.
\end{attention}

\etape{Confirmer ou refuser}

Deux droits distincts existent : \textbf{confirmer} un retrait et le
\textbf{refuser}. Refusez si le code ne correspond pas.

\begin{resultat}
La commande passe à \textbf{Retirée} et le colis à \textbf{Remis}.
\end{resultat}

\subsection{Les expéditions}

Une expédition suit ce cycle :

\begin{illustration}
\begin{center}
\begin{tikzpicture}[
  node distance=4mm,
  et/.style={rectangle,rounded corners=3pt,line width=0.8pt,
             minimum height=9mm,text width=22mm,align=center,
             font=\scriptsize\bfseries},
  f/.style={-{Stealth[length=2.6mm]},garahAttenue,line width=0.9pt}
]
\node[et,draw=garahMarque,fill=garahMarque!8] (a) {Créée};
\node[et,draw=garahMarque,fill=garahMarque!8,right=of a] (b) {Préparée};
\node[et,draw=garahMarque,fill=garahMarque!8,right=of b] (c) {En transit};
\node[et,draw=garahSucces,fill=garahSucces!10,right=of c] (d) {Disponible};
\node[et,draw=garahSucces,fill=garahSucces!10,right=of d] (e) {Remise};
\node[et,draw=garahAlerte,fill=garahAlerte!12,below=6mm of c] (x) {Bloquée};
\draw[f] (a) -- (b); \draw[f] (b) -- (c); \draw[f] (c) -- (d); \draw[f] (d) -- (e);
\draw[f] (c) -- (x);
\end{tikzpicture}
\end{center}
\end{illustration}
\legendeIllus{3.4}{Le cycle d'une expédition. « Bloquée » signale un incident :
il doit être traité, sinon le client attend sans savoir.}

\begin{astuce}
\textbf{Lieux} et \textbf{Itinéraires} sont des référentiels : on y vient
rarement, pour ajouter un point ou un chemin. \textbf{Expéditions} et
\textbf{Retraits} sont le travail quotidien. C'est pourquoi le menu les place
dans cet ordre.
\end{astuce}

% -----------------------------------------------------------------------------
\section{Module 4 — Le catalogue et le stock}

\subsection{Tutoriel : ajouter un produit}

\begin{objectif}
Créer un article vendable, avec ses déclinaisons et ses paliers de prix.
\end{objectif}

\begin{avant}
Vous devez avoir le droit de créer un produit. Le \textbf{marchand} et la
\textbf{catégorie} doivent déjà exister.
\end{avant}

\etape{Ouvrir les produits}

Menu \textbf{Le catalogue $\rightarrow$ Produits}, puis le bouton de création.

\etape{Renseigner l'article}

Nom, description, catégorie, marchand, photos.

\begin{astuce}
Mettez \textbf{plusieurs photos}. Une seule photo d'un article qu'on ne peut ni
toucher ni essayer est une raison de ne pas acheter. La première photo est
celle qui part au panier.
\end{astuce}

\etape{Créer les déclinaisons}

Une déclinaison = une version achetable. C'est elle qui porte le stock et le
prix, pas l'article.

\etape{Saisir les paliers de prix}

\begin{attention}
\textbf{Le dernier palier doit être ouvert} — c'est-à-dire sans quantité
maximale, ce qui signifie « et au-delà ».

Si vous fermez le dernier palier à 10, un client qui demande 11 tombe dans un
trou : \emph{aucun prix n'existe pour cette quantité}. La boutique refuse alors
la vente et affiche « Prix non fixé pour cette quantité ».
\end{attention}

\begin{illustration}
\begin{center}
\begin{tikzpicture}[x=1mm,y=1mm]
  % Bon
  \node[font=\scriptsize\bfseries,text=garahSucces,anchor=west] at (0,26) {\ding{51} Grille correcte};
  \draw[draw=garahSucces,fill=garahSucces!7,rounded corners=2pt] (0,14) rectangle (28,22);
  \node[font=\tiny] at (14,18) {1 -- 6};
  \draw[draw=garahSucces,fill=garahSucces!7,rounded corners=2pt] (31,14) rectangle (59,22);
  \node[font=\tiny] at (45,18) {7 -- 10};
  \draw[draw=garahSucces,fill=garahSucces!7,rounded corners=2pt] (62,14) rectangle (90,22);
  \node[font=\tiny\bfseries] at (76,18) {11 et +};
  \node[font=\tiny,text=garahAttenue,anchor=west] at (0,9) {Aucune quantité sans prix.};

  % Mauvais
  \node[font=\scriptsize\bfseries,text=garahDanger,anchor=west] at (100,26) {\ding{55} Grille fermée};
  \draw[draw=garahDanger,fill=garahDanger!7,rounded corners=2pt] (100,14) rectangle (124,22);
  \node[font=\tiny] at (112,18) {1 -- 6};
  \draw[draw=garahDanger,fill=garahDanger!7,rounded corners=2pt] (127,14) rectangle (151,22);
  \node[font=\tiny] at (139,18) {7 -- 10};
  \node[font=\tiny,text=garahDanger,anchor=west] at (100,9) {À partir de 11 : plus aucun prix.};
\end{tikzpicture}
\end{center}
\end{illustration}
\legendeIllus{3.5}{À gauche, une grille dont le dernier palier est ouvert :
toute quantité a un prix. À droite, une grille fermée à 10 : au-delà, la vente
est impossible.}

\etape{Publier}

Un produit est \textbf{Brouillon}, \textbf{Publié} ou \textbf{Masqué}. Seul un
produit publié apparaît dans la boutique.

\begin{resultat}
L'article apparaît dans le catalogue de la boutique, dans sa catégorie, avec sa
grille de prix.
\end{resultat}

\subsection{Le stock}

\begin{attention}
Le stock se décide \textbf{par déclinaison}. Une déclinaison à zéro s'affiche
« Épuisée » et barrée dans la boutique — elle reste \emph{visible}, car la
masquer ferait croire qu'elle n'existe pas et enverrait chercher ailleurs.
\end{attention}

% -----------------------------------------------------------------------------
\section{Module 5 — L'après-vente}

\subsection{Traiter un retour}

Ouvrez \textbf{Après-vente $\rightarrow$ Retours}. Le circuit
(figure~2.9) est :
\textbf{Demandé} $\rightarrow$ \textbf{Accepté} ou \textbf{Refusé}
$\rightarrow$ \textbf{Réceptionné} $\rightarrow$ \textbf{Validé}.

Chaque passage est une décision qui vous engage :

\begin{center}
\begin{tabular}{@{}>{\bfseries}p{3.4cm} p{11.4cm}@{}}
\toprule
Étape & Ce que vous décidez \\
\midrule
Accepter & Vous vous engagez à reprendre la marchandise. \\
Refuser & Vous refusez la reprise. Expliquez pourquoi au client. \\
Réceptionner & Vous confirmez avoir \emph{physiquement} reçu l'article. \\
Valider & Vous clôturez le dossier. \\
\bottomrule
\end{tabular}
\end{center}

\subsection{Traiter une réclamation}

Ouvrez \textbf{Après-vente $\rightarrow$ Réclamations}. Une réclamation passe de
\textbf{Ouverte} à \textbf{En cours} lorsque vous vous en saisissez, puis à
\textbf{Résolue}.

% -----------------------------------------------------------------------------
\section{Module 6 — Travailler avec ses collègues}

\subsection{La messagerie interne}

Le menu \textbf{Messages} permet d'écrire à un collègue. C'est un canal
\textbf{interne} : les clients n'y ont pas accès.

\begin{attention}
\textbf{Ne confondez pas la messagerie interne et le service client.} Le service
client parle \emph{aux clients}. La messagerie parle \emph{entre collègues}.
Écrire au mauvais endroit expose une discussion interne, ou laisse un client
sans réponse.
\end{attention}

\subsection{Mon service}

\textbf{Les gens $\rightarrow$ Mon service} montre les membres de votre service.
C'est utile pour savoir à qui s'adresser.

% -----------------------------------------------------------------------------
\section{Scénario complet : traiter une nouvelle demande client}

\begin{tcolorbox}[colback=garahFond,colframe=garahEncre,arc=4pt,boxrule=1pt]
\textbf{Situation.} Un client vient de poser une question sur un article. La
conversation est en attente depuis trois minutes.
\end{tcolorbox}

\begin{enumerate}[leftmargin=*,itemsep=4pt]
  \item Vous vous connectez au back-office.
  \item Le \textbf{tableau de bord} vous signale des conversations en attente.
  \item Vous ouvrez \textbf{Après-vente $\rightarrow$ Service client}, onglet
        \textbf{En attente}.
  \item Vous repérez le dossier : sujet, nom du client, code client, et le
        compteur de messages non lus.
  \item Vous cliquez sur \textbf{Prendre}. Le dossier passe à \textbf{En cours}
        et vous est affecté.
  \item Vous ouvrez le fil et lisez la question.
  \item Si besoin, vous consultez la fiche de l'article dans un autre onglet
        pour vérifier le stock et les paliers.
  \item Vous répondez. \textbf{Le client vous voit répondre en direct.}
  \item S'il négocie, vous lui faites une \textbf{proposition de prix} — pensez
        à sa date d'expiration.
  \item Une fois la demande réglée, vous \textbf{clôturez} la conversation.
\end{enumerate}

\begin{exercice}
\textbf{Objectif.} Traiter une conversation de bout en bout.

\medskip
\textbf{Actions attendues}
\begin{enumerate}[leftmargin=*,itemsep=1pt,topsep=2pt]
  \item Depuis un second appareil, ouvrir une conversation en tant que client.
  \item Côté back-office, la retrouver dans l'onglet « En attente ».
  \item La prendre en charge.
  \item Répondre, et \emph{vérifier que la réponse apparaît côté client sans
        recharger}.
  \item Faire répondre le client, et vérifier l'arrivée en direct côté
        back-office.
  \item Clore la conversation.
\end{enumerate}

\medskip
\textbf{Résultat attendu.} Le dossier est « Fermée », le compteur de non lus est
à zéro, et le client ne peut plus écrire dedans.

\medskip
\textbf{Critères de réussite}
\begin{itemize}[leftmargin=*,itemsep=1pt,topsep=2pt]
  \item les messages sont arrivés \emph{sans rechargement} dans les deux sens ;
  \item le compteur n'a \emph{pas} augmenté quand vous avez répondu ;
  \item vous savez expliquer pourquoi rien n'arrivait avant la prise en charge.
\end{itemize}
\end{exercice}

\begin{exercice}
\textbf{Exercice avancé — la grille de prix piégée}

\medskip
Créez un article dont le dernier palier est \textbf{fermé} à 10 pièces. Puis,
côté boutique, essayez d'en commander 11.

\medskip
\textbf{Résultat attendu.} La boutique refuse et affiche « Prix non fixé pour
cette quantité » ; le sélecteur de quantité s'arrête à 10 et explique pourquoi.

\medskip
\textbf{Puis corrigez} la grille en ouvrant le dernier palier, et vérifiez que
la commande de 11 devient possible.

\medskip
\textbf{Critère de réussite.} Vous savez dire, sans regarder ce manuel, ce
qu'une grille fermée provoque côté client.
\end{exercice}

\begin{checklist}
\case Je sais me connecter au back-office\\
\case Je sais expliquer pourquoi un écran n'apparaît pas dans mon menu\\
\case Je sais retrouver un écran à partir de la question qu'il traite\\
\case Je sais filtrer la file du service client\\
\case Je sais prendre en charge une conversation\\
\case Je sais ce que fait « Clore », et que c'est irréversible pour le client\\
\case Je sais lire le compteur de messages non lus\\
\case Je sais reconnaître l'état d'une commande et dire ce que j'en fais\\
\case Je sais que je ne prépare jamais une commande non payée\\
\case Je sais confirmer un retrait, et pourquoi le code est indispensable\\
\case Je sais suivre une expédition et repérer un colis bloqué\\
\case Je sais créer un produit, ses déclinaisons et ses paliers\\
\case Je sais pourquoi le dernier palier doit rester ouvert\\
\case Je sais traiter un retour et une réclamation, et les distinguer\\
\case Je sais écrire à un collègue sans me tromper de canal
\end{checklist}

% =============================================================================
\chapter{Formation de l'Administrateur}
\label{ch:admin}
% =============================================================================

\begin{roleBox}
\textbf{Qui vous êtes.} Vous ne faites pas le travail quotidien : vous
\textbf{décidez qui le fait}. Vous créez les comptes de l'équipe, vous
attribuez les profils, vous gérez les marchands et vous suivez l'argent dû.

\medskip
\textbf{Ce que vous pouvez faire, en plus d'un responsable}
\begin{itemize}[leftmargin=*,itemsep=2pt,topsep=3pt]
  \item créer, modifier et désactiver des comptes internes ;
  \item créer des profils et choisir les droits qu'ils contiennent ;
  \item attribuer un ou plusieurs profils à un responsable ;
  \item gérer les marchands et suivre ce qui leur est dû ;
  \item annuler une commande déjà payée, avec remboursement ;
  \item consulter les statistiques générales.
\end{itemize}

\medskip
\textbf{Ce que vous ne pouvez pas faire}
\begin{itemize}[leftmargin=*,itemsep=2pt,topsep=3pt]
  \item inventer un droit qui n'existe pas dans le catalogue des permissions ;
  \item désactiver le compte d'amorçage — plus personne ne pourrait rendre la
        main ;
  \item consulter le journal des actions, réservé au super administrateur.
\end{itemize}
\end{roleBox}

% -----------------------------------------------------------------------------
\section{Module 1 — Comprendre le système des droits}

C'est le concept le plus important de votre rôle. Il tient en une chaîne :

\begin{illustration}
\begin{center}
\begin{tikzpicture}[
  b/.style={rectangle,rounded corners=4pt,line width=1pt,
            minimum height=15mm,text width=34mm,align=center,font=\small},
  f/.style={-{Stealth[length=3.4mm]},garahAttenue,line width=1.2pt}
]
\node[b,draw=garahAccent,fill=garahAccent!8] (a)
  {\textbf{ADMINISTRATEUR}\\[2pt]\scriptsize décide qui fait quoi};
\node[b,draw=garahPrimaire,fill=garahPrimaire!8,right=18mm of a] (b)
  {\textbf{PROFIL}\\[2pt]\scriptsize un paquet de droits};
\node[b,draw=garahMarque,fill=garahMarque!8,right=18mm of b] (c)
  {\textbf{RESPONSABLE}\\[2pt]\scriptsize fait le travail};
\draw[f] (a) -- (b);
\draw[f] (b) -- (c);
\end{tikzpicture}
\end{center}
\end{illustration}
\legendeIllus{4.1}{La chaîne des droits, telle qu'elle est rappelée en haut de
l'écran « Équipe ».}

\begin{retenir}
\textbf{Un compte n'a aucun droit en propre.} Tous ses droits viennent de ses
profils. Retirer un profil retire les droits correspondants, immédiatement.
\end{retenir}

\begin{attention}
\textbf{Un profil ne CRÉE pas de droits : il en SÉLECTIONNE.} Vous choisissez
parmi un catalogue de permissions existantes. C'est ce qui vous empêche
d'inventer un droit qui ne correspondrait à rien — et qui ne garderait donc
rien.
\end{attention}

\subsection{Tutoriel : créer un compte et lui donner des droits}

\begin{objectif}
Créer un compte interne et le rendre opérationnel.
\end{objectif}

\etape{Ouvrir l'équipe}

Menu \textbf{Les gens $\rightarrow$ Équipe}.

\begin{illustration}
\begin{center}
\begin{tikzpicture}[x=1mm,y=1mm]
  % Entête
  \node[font=\small\bfseries,anchor=west] at (0,50) {Équipe};
  \node[font=\tiny,text=garahAttenue,anchor=west] at (0,45) {3 compte(s) interne(s)};
  \draw[fill=garahSucces,draw=none,rounded corners=2pt] (96,44) rectangle (120,52);
  \node[font=\tiny,text=white] at (108,48) {Profils};
  \draw[fill=garahPrimaire,draw=none,rounded corners=2pt] (123,44) rectangle (152,52);
  \node[font=\tiny,text=white] at (137.5,48) {Nouveau compte};
  \node at (108,57) {\num{2}};
  \node at (137.5,57) {\num{1}};

  % Bandeau explicatif
  \draw[draw=garahBordure,fill=garahFond,rounded corners=3pt] (0,28) rectangle (152,40);
  \node[font=\tiny\bfseries,anchor=west] at (4,36) {Comment se composent les accès};
  \node[font=\tiny,text=garahAttenue,anchor=west] at (4,31.5)
    {Un compte n'a pas de droits en propre : il en reçoit par ses profils.};

  % Ligne de tableau
  \draw[draw=garahBordure,fill=white,rounded corners=3pt] (0,4) rectangle (152,24);
  \node[font=\tiny\bfseries,anchor=west] at (4,18) {Ophelie BRIA};
  \node[font=\tiny,text=garahAttenue,anchor=west] at (4,13) {bria@gmail.com};
  \draw[fill=garahAlerte!18,draw=none,rounded corners=5pt] (48,15) rectangle (76,21);
  \node[font=\tiny] at (62,18) {RESPONSABLE};
  \draw[fill=garahMarque!18,draw=none,rounded corners=5pt] (82,15) rectangle (116,21);
  \node[font=\tiny] at (99,18) {SERVICE MARKETING};
  \node at (99,10) {\num{3}};
\end{tikzpicture}
\end{center}
\end{illustration}
\legendeIllus{4.2}{L'écran Équipe. \num{1} créer un compte. \num{2} gérer les
profils et leurs droits. \num{3} les profils attribués à ce compte — c'est
là que se lisent ses accès réels.}

\etape{Créer le compte}

Cliquez sur \textbf{« Nouveau compte »} et renseignez nom, prénom, adresse
e-mail.

\begin{attention}
Un compte créé mais \textbf{sans profil ne peut rien faire}. Il se connectera
et verra un menu presque vide. Ne vous arrêtez pas là.
\end{attention}

\etape{Attribuer un ou plusieurs profils}

Choisissez les profils qui correspondent au métier de la personne.

\begin{astuce}
Donnez \textbf{le moins de droits possible}, quitte à en ajouter ensuite. Un
droit accordé « au cas où » finit par servir un jour où il ne fallait pas.
\end{astuce}

\etape{Vérifier}

La ligne du tableau doit afficher le niveau et les profils. Demandez à la
personne de se connecter et de vérifier son menu.

\begin{resultat}
Le compte apparaît \textbf{Actif}, avec ses profils. La personne voit dans son
menu exactement les écrans que ses profils autorisent.
\end{resultat}

\subsection{Que faire si ça ne fonctionne pas ?}

\begin{longtable}{@{}p{4.2cm} p{4cm} p{6cm}@{}}
\toprule
\bfseries Problème & \bfseries Cause possible & \bfseries Solution \\
\midrule
\endhead
La personne ne voit pas un écran
  & Aucun profil ne porte le droit correspondant
  & Ajoutez le droit au profil, ou attribuez un autre profil \\[4pt]
Elle ne peut pas se connecter
  & Adresse non confirmée, ou compte inactif
  & Vérifiez la mention « NON VÉRIFIÉE » et le statut \\[4pt]
Le profil est refusé
  & Le profil est désactivé
  & Réactivez-le, ou choisissez-en un autre \\[4pt]
Impossible de désactiver un compte
  & C'est le compte d'amorçage
  & C'est volontaire : le désactiver ferait perdre la main à tout le monde \\
\bottomrule
\end{longtable}

% -----------------------------------------------------------------------------
\section{Module 2 — Les marchands et l'argent}

\subsection{Les marchands}

Menu \textbf{Le catalogue $\rightarrow$ Marchands}. Un marchand fournit des
produits.

\begin{astuce}
Le marchand est rangé dans « Le catalogue » et non dans « Les gens », parce
qu'on le choisit \textbf{en créant un produit}, pas en gérant des comptes.
\end{astuce}

\subsection{Marchands à payer}

Menu \textbf{Les ventes $\rightarrow$ Marchands à payer}. C'est l'écran des
dettes envers les marchands.

\begin{retenir}
Une vente produit \textbf{trois choses} : une commande, un encaissement, et une
dette envers le marchand. C'est pourquoi les trois écrans se suivent dans le
menu.
\end{retenir}

\subsection{Les statistiques}

Menu \textbf{Pilotage $\rightarrow$ Statistiques}.

\captureAInserer[9em]{L'écran Statistiques de votre installation.}

% -----------------------------------------------------------------------------
\section{Module 3 — Les situations particulières}

\subsection{Annuler une commande payée}

\begin{attention}
\textbf{Vous êtes le seul à pouvoir le faire}, et cela déclenche un
remboursement. Ce n'est pas une correction de saisie : c'est une décision
financière qui laisse une trace dans le journal des actions.

Vérifiez trois fois avant de valider.
\end{attention}

\subsection{Retirer un dossier à quelqu'un}

Une conversation prise par un collègue absent peut être remise en attente.
L'opération est \textbf{tracée} : retirer un dossier à quelqu'un est une
décision qui doit pouvoir s'expliquer.

\begin{exercice}
\textbf{Objectif.} Créer un responsable opérationnel, sans lui donner trop.

\medskip
\textbf{Actions attendues}
\begin{enumerate}[leftmargin=*,itemsep=1pt,topsep=2pt]
  \item Créer un profil qui ne contient que les droits du service client.
  \item Créer un compte et lui attribuer ce seul profil.
  \item Faire connecter la personne et \emph{lister avec elle} les écrans
        visibles.
  \item Vérifier qu'elle ne voit ni « Équipe », ni « Marchands à payer », ni
        « Journal des actions ».
\end{enumerate}

\medskip
\textbf{Critères de réussite.} Le menu de la personne contient le service
client et rien de plus. Elle sait expliquer pourquoi.
\end{exercice}

\begin{checklist}
\case Je sais expliquer la chaîne administrateur $\rightarrow$ profil $\rightarrow$ responsable\\
\case Je sais créer un compte interne\\
\case Je sais créer un profil et y choisir des droits\\
\case Je sais attribuer et retirer un profil\\
\case Je sais vérifier les accès réels d'une personne\\
\case Je sais pourquoi un profil ne peut pas inventer un droit\\
\case Je sais gérer les marchands\\
\case Je sais lire l'écran « Marchands à payer »\\
\case Je sais qu'annuler une commande payée est une décision financière tracée\\
\case Je sais pourquoi le compte d'amorçage ne se désactive pas
\end{checklist}

% =============================================================================
\chapter{Formation du Super administrateur}
\label{ch:superadmin}
% =============================================================================

\begin{roleBox}
\textbf{Qui vous êtes.} Vous avez tous les droits, et une responsabilité que
les autres n'ont pas : \textbf{savoir ce qui s'est passé}.

\medskip
\textbf{Ce que vous seul pouvez faire}
\begin{itemize}[leftmargin=*,itemsep=2pt,topsep=3pt]
  \item consulter le \textbf{journal des actions} internes ;
  \item accéder aux écrans de surveillance et de sécurité.
\end{itemize}

\begin{attention}
\textbf{Ce que vous voyez dans le journal est sensible.} Les clichés « avant /
après » contiennent l'état des objets modifiés — potentiellement n'importe
quelle donnée du système. C'est pour cette raison que la route est réservée à
votre seul rôle.
\end{attention}
\end{roleBox}

% -----------------------------------------------------------------------------
\section{Module 1 — Le journal des actions}

\subsection{À quelle question il répond}

\begin{center}
\begin{tcolorbox}[width=0.9\linewidth,colback=garahAccent!7,
  colframe=garahAccent,arc=4pt,boxrule=1pt]
\centering
« Qui a annulé cette commande ? »\quad « Qui a mis ce taux à 30\,\% ? »\\
« Qui a remboursé ce client ? »
\end{tcolorbox}
\end{center}

Ces questions se posent toutes \textbf{des semaines après le geste}, quand la
table concernée ne porte plus que son état courant.

\begin{retenir}
\textbf{Il y a trois journaux, et celui-ci n'est pas les deux autres :}
\begin{itemize}[leftmargin=*,itemsep=1pt,topsep=3pt]
  \item ce que fait un \textbf{client} $\rightarrow$ son parcours, mesuré à part ;
  \item les faits d'\textbf{authentification} $\rightarrow$ le journal de sécurité ;
  \item les actions \textbf{internes} $\rightarrow$ le journal des actions.
\end{itemize}
Les fondre en un seul les rendrait tous les trois inexploitables : les quelques
gestes internes d'une journée disparaîtraient sous le trafic de la boutique.
\end{retenir}

\subsection{Tutoriel : retrouver qui a fait quoi}

\begin{objectif}
Retrouver une action précise et lire exactement ce qui a changé.
\end{objectif}

\etape{Ouvrir le journal}

Menu \textbf{La trace $\rightarrow$ Journal des actions}.

\etape{Filtrer}

Deux listes : \textbf{le geste} et \textbf{l'objet}. Elles sont peuplées par ce
que le journal \emph{contient} — un filtre proposé rend donc toujours quelque
chose.

\etape{Déplier une ligne}

Cliquez sur une ligne. Un tableau montre \textbf{uniquement les champs qui ont
changé} : l'ancienne valeur barrée, la nouvelle mise en avant.

\begin{illustration}
\begin{center}
\begin{tikzpicture}[x=1mm,y=1mm]
  \draw[fill=garahSurfaceNuit,draw=garahAccent!30,rounded corners=3pt] (0,0) rectangle (150,40);
  \node[font=\tiny\bfseries,text=garahAttenue,anchor=west] at (5,35) {CE QUI A CHANGÉ};
  \node[font=\tiny\bfseries,text=garahAttenue,anchor=west] at (62,35) {AVANT};
  \node[font=\tiny\bfseries,text=garahAttenue,anchor=west] at (105,35) {APRÈS};
  \draw[garahAccent!25] (5,31) -- (145,31);

  \node[font=\tiny\bfseries,text=white,anchor=west] at (5,25) {Prénom};
  \node[font=\tiny,text=garahAttenue,anchor=west] at (62,25) {Ophelie};
  \draw[garahAttenue] (62,25) -- (74,25);
  \draw[garahMarque,line width=1.2pt] (103,21) -- (103,29);
  \node[font=\tiny\bfseries,text=white,anchor=west] at (105,25) {Ophélie};
  \draw[garahAccent!25] (5,19) -- (145,19);

  \node[font=\tiny\bfseries,text=white,anchor=west] at (5,13) {Statut};
  \node[font=\tiny,text=garahAttenue,anchor=west] at (62,13) {ACTIF};
  \draw[garahAttenue] (62,13) -- (72,13);
  \draw[garahMarque,line width=1.2pt] (103,9) -- (103,17);
  \node[font=\tiny\bfseries,text=white,anchor=west] at (105,13) {INACTIF};

  \node at (156,25) {\num{1}};
  \node at (156,13) {\num{2}};
\end{tikzpicture}
\end{center}
\end{illustration}
\legendeIllus{5.1}{Le détail d'une action. \num{1} l'ancienne valeur, barrée et
atténuée. \num{2} la nouvelle, portée par un liseré. Seuls les champs
\emph{modifiés} sont affichés.}

\begin{astuce}
Si l'écran affiche \textbf{« Aucun champ n'a changé »}, ce n'est pas une panne :
le geste a bien été enregistré, mais les valeurs sont restées les mêmes. C'est
une information en soi.
\end{astuce}

\begin{attention}
Les \textbf{identifiants restent des identifiants} : le journal écrit « n\textdegree~3 »
et non le nom du profil. C'est volontaire — afficher le nom d'\emph{aujourd'hui}
sur un geste d'il y a six mois serait un mensonge. Un journal d'audit montre ce
qui a été enregistré, pas ce qui est vrai maintenant.
\end{attention}

\begin{resultat}
Vous savez qui a agi, quand, sur quel objet, depuis quelle adresse, et
exactement ce qui a changé.
\end{resultat}

\begin{checklist}
\case Je sais ouvrir le journal des actions\\
\case Je sais filtrer par geste et par objet\\
\case Je sais lire un détail « avant / après »\\
\case Je sais interpréter « Aucun champ n'a changé »\\
\case Je sais pourquoi le journal montre des identifiants et non des noms\\
\case Je sais distinguer les trois journaux\\
\case Je sais que le contenu du journal est sensible et pourquoi il m'est réservé
\end{checklist}

% =============================================================================
\chapter{Dépannage général}
\label{ch:depannage}
% =============================================================================

Ce chapitre rassemble les problèmes qui touchent \textbf{tous les acteurs}. La
méthode est toujours la même :

\begin{illustration}
\begin{center}
\begin{tikzpicture}[
  b/.style={rectangle,rounded corners=3pt,draw=garahMarque,fill=garahMarque!8,
            line width=0.9pt,minimum height=11mm,text width=26mm,align=center,
            font=\small\bfseries},
  f/.style={-{Stealth[length=3mm]},garahAttenue,line width=1pt}
]
\node[b] (a) {Problème};
\node[b,right=8mm of a] (b) {Cause possible};
\node[b,right=8mm of b] (c) {Vérification};
\node[b,right=8mm of c] (d) {Solution};
\draw[f] (a) -- (b); \draw[f] (b) -- (c); \draw[f] (c) -- (d);
\end{tikzpicture}
\end{center}
\end{illustration}
\legendeIllus{6.1}{La méthode de dépannage. Ne sautez pas la vérification :
c'est elle qui distingue une vraie panne d'un malentendu.}

\section{Connexion et accès}

\begin{longtable}{@{}p{3.6cm} p{3.4cm} p{3.4cm} p{3.8cm}@{}}
\toprule
\bfseries Problème & \bfseries Cause & \bfseries Vérification & \bfseries Solution \\
\midrule
\endhead
« Mauvais identifiants »
  & Adresse ou mot de passe faux
  & Retapez lentement
  & Le message ne dit pas lequel des deux est faux : le préciser aiderait
    quelqu'un qui cherche à deviner \\[4pt]
Déconnexion inattendue
  & La session a expiré
  & Reconnectez-vous
  & Normal. L'application mobile garde la session plus longtemps \\[4pt]
« Accès refusé »
  & Vos profils ne portent pas ce droit
  & Regardez si l'entrée est dans votre menu
  & Demandez le droit à un administrateur \\[4pt]
Un écran manque au menu
  & Même cause
  & Comparez avec un collègue du même métier
  & Ce n'est pas une panne : le menu ne montre que l'accessible \\[4pt]
« NON VÉRIFIÉE » à côté du compte
  & Adresse e-mail non confirmée
  & Cherchez le message de confirmation
  & Cliquez le lien reçu. Vérifiez les indésirables \\
\bottomrule
\end{longtable}

\section{Commande et paiement}

\begin{longtable}{@{}p{3.6cm} p{3.4cm} p{3.4cm} p{3.8cm}@{}}
\toprule
\bfseries Problème & \bfseries Cause & \bfseries Vérification & \bfseries Solution \\
\midrule
\endhead
« Votre panier est vide » alors qu'il ne l'est pas
  & Le panier n'a pas été envoyé au serveur
  & Revenez au panier : les articles y sont-ils ?
  & Repassez par « Passer commande ». Si le réseau a échoué, un message le dit
    et le panier est conservé \\[4pt]
« Le numéro de téléphone est obligatoire »
  & Le champ Mobile Money est vide
  & Regardez le champ au-dessus du bouton
  & Saisissez le numéro qui recevra la demande de validation \\[4pt]
Le paiement reste « en attente »
  & L'opérateur n'a pas encore confirmé
  & Avez-vous validé sur votre téléphone ?
  & Validez, puis appuyez sur le bouton qui redemande l'état. L'écran ne se
    rafraîchit pas seul, exprès \\[4pt]
« Prix non fixé pour cette quantité »
  & Aucun palier ne couvre cette quantité
  & Regardez la grille de prix
  & Baissez la quantité, ou contactez le marchand. Côté back-office : ouvrez
    le dernier palier \\[4pt]
Le bouton « + » ne répond plus
  & Le stock ou le tarif borne
  & Lisez le message sous le compteur
  & Les deux cas n'appellent pas le même geste \\[4pt]
Total différent du sous-total
  & Les frais d'acheminement s'ajoutent
  & Changez de point de récupération
  & C'est normal : les frais dépendent du point \\
\bottomrule
\end{longtable}

\section{Conversations et temps réel}

\begin{longtable}{@{}p{3.6cm} p{3.4cm} p{3.4cm} p{3.8cm}@{}}
\toprule
\bfseries Problème & \bfseries Cause & \bfseries Vérification & \bfseries Solution \\
\midrule
\endhead
Les messages n'arrivent pas en direct
  & La conversation n'est pas affectée
  & L'étiquette dit-elle « En attente » ?
  & Prenez-la en charge. Sans responsable, rien n'est poussé \\[4pt]
Rien n'arrive même après la prise en charge
  & La connexion temps réel est coupée
  & Rechargez la page une fois
  & Si le rechargement affiche les messages, signalez-le : le direct est en
    cause, pas les données \\[4pt]
Impossible d'écrire
  & La conversation est fermée
  & Regardez l'étiquette
  & Ouvrez-en une nouvelle \\[4pt]
« Prendre » échoue
  & Un collègue a été plus rapide
  & Actualisez la liste
  & C'est ce qui empêche deux réponses simultanées \\
\bottomrule
\end{longtable}

\section{Affichage et réseau}

\begin{longtable}{@{}p{3.6cm} p{3.4cm} p{3.4cm} p{3.8cm}@{}}
\toprule
\bfseries Problème & \bfseries Cause & \bfseries Vérification & \bfseries Solution \\
\midrule
\endhead
Le catalogue paraît presque vide
  & Un filtre est actif
  & Cherchez la puce de filtre
  & Enlevez-la avec sa croix \\[4pt]
Une image ne s'affiche pas
  & L'adresse signée a expiré
  & Rechargez l'écran
  & Une silhouette s'affiche à la place plutôt qu'une icône cassée \\[4pt]
« Pas de connexion »
  & Réseau absent ou coupé
  & Ouvrez une autre page
  & Réessayez. Le panier local, lui, n'est pas perdu \\
\bottomrule
\end{longtable}

% =============================================================================
\chapter{Guide rapide}
\label{ch:rapide}
% =============================================================================

Cette partie s'adresse à quelqu'un \textbf{qui sait déjà}, et qui cherche
seulement à retrouver un chemin.

\section{Où trouver quoi, dans le back-office}

\begin{longtable}{@{}>{\bfseries}p{5.6cm} p{9.2cm}@{}}
\toprule
Je veux\dots & Menu \\
\midrule
\endhead
voir où on en est & Pilotage $\rightarrow$ Tableau de bord \\
voir les chiffres & Pilotage $\rightarrow$ Statistiques \\
ajouter un produit & Le catalogue $\rightarrow$ Produits \\
ranger les catégories & Le catalogue $\rightarrow$ Catégories \\
gérer les dimensions de déclinaison & Le catalogue $\rightarrow$ Dimensions \\
vérifier un stock & Le catalogue $\rightarrow$ Stock \\
gérer un fournisseur & Le catalogue $\rightarrow$ Marchands \\
traiter une commande & Les ventes $\rightarrow$ Commandes \\
vérifier un encaissement & Les ventes $\rightarrow$ Paiements \\
savoir ce qu'on doit aux marchands & Les ventes $\rightarrow$ Marchands à payer \\
suivre un colis & L'acheminement $\rightarrow$ Expéditions \\
remettre la marchandise & L'acheminement $\rightarrow$ Retraits \\
ajouter un point de récupération & L'acheminement $\rightarrow$ Lieux \\
définir un chemin & L'acheminement $\rightarrow$ Itinéraires \\
traiter un problème client & Après-vente $\rightarrow$ Réclamations \\
traiter une marchandise qui revient & Après-vente $\rightarrow$ Retours \\
répondre à un client & Après-vente $\rightarrow$ Service client \\
consulter un client & Les gens $\rightarrow$ Clients \\
gérer l'équipe & Les gens $\rightarrow$ Équipe \\
voir mes collègues & Les gens $\rightarrow$ Mon service \\
écrire à un collègue & Messages \\
savoir qui a fait quoi & La trace $\rightarrow$ Journal des actions \\
\bottomrule
\end{longtable}

\section{Les gestes du client, en une page}

\begin{longtable}{@{}>{\bfseries}p{5.6cm} p{9.2cm}@{}}
\toprule
Je veux\dots & Où \\
\midrule
\endhead
parcourir sans compte & Catalogue, ou une puce de catégorie sur l'accueil \\
mettre de côté & le cœur sur une fiche ($\rightarrow$ Ma liste) \\
commander & Panier $\rightarrow$ Passer commande \\
payer & Choisir le point, puis Mobile Money \\
suivre un colis \emph{sans compte} & Suivre un colis + numéro de suivi \\
retirer & se présenter avec le \textbf{code de retrait} \\
poser une question & « Contacter » sur une fiche article \\
voir mes échanges & Mes discussions \\
voir mes achats & Mes commandes \\
renvoyer un article & Mes retours \\
signaler un problème & Mes réclamations \\
\bottomrule
\end{longtable}

\section{Les cinq réflexes qui évitent les erreurs}

\begin{enumerate}[leftmargin=*,itemsep=5pt]
  \item \textbf{Le dernier palier de prix reste toujours ouvert.} Sinon la
        vente devient impossible au-delà.
  \item \textbf{On ne prépare jamais une commande non payée.}
  \item \textbf{On ne confirme jamais un retrait sans code.}
  \item \textbf{On ne clôture une conversation que si le sujet est réglé} : le
        client ne pourra plus y écrire.
  \item \textbf{On donne le moins de droits possible}, quitte à en ajouter.
\end{enumerate}

% =============================================================================
\chapter{Manuel du formateur}
\label{ch:formateur}
% =============================================================================

\section{Avant la session}

\begin{center}
\begin{tabular}{@{}>{\bfseries}p{4.4cm} p{10.4cm}@{}}
\toprule
À préparer & Pourquoi \\
\midrule
Un environnement d'essai & Personne n'apprend en touchant les vraies commandes. \\
Un compte par stagiaire & Se connecter fait partie de la formation. \\
Deux appareils par binôme & Le temps réel ne se démontre qu'à deux écrans. \\
Un jeu de données & Au moins deux catégories, un article à plusieurs paliers,
                    un article épuisé. \\
Les captures de \emph{votre} installation & À insérer aux emplacements prévus. \\
\bottomrule
\end{tabular}
\end{center}

\begin{attention}
\textbf{Préparez un article dont le dernier palier est fermé.} C'est le meilleur
exercice du manuel : il montre en trente secondes pourquoi une règle de saisie
compte, et l'erreur se voit côté client.
\end{attention}

\section{Programme proposé — deux journées}

\subsection*{Journée 1 — Comprendre et vendre}

\begin{center}
\begin{tabular}{@{}>{\bfseries}p{2.2cm} p{5.4cm} p{6.8cm}@{}}
\toprule
Durée & Séquence & Démonstration à faire \\
\midrule
45 min & Découvrir GARAH (ch.~\ref{ch:decouvrir}) & Le parcours d'une commande au tableau \\
30 min & Vocabulaire & Montrer une déclinaison et une grille de paliers \\
1 h 30 & Le client (ch.~\ref{ch:client}) & Commander en direct, du catalogue au paiement \\
30 min & Exercice client & Chacun passe une commande \\
1 h & Suivi et retrait & Faire circuler un code de retrait \\
1 h & Conversations & \textbf{Deux écrans côte à côte}, temps réel \\
\bottomrule
\end{tabular}
\end{center}

\subsection*{Journée 2 — Traiter et administrer}

\begin{center}
\begin{tabular}{@{}>{\bfseries}p{2.2cm} p{5.4cm} p{6.8cm}@{}}
\toprule
Durée & Séquence & Démonstration à faire \\
\midrule
30 min & Le menu par question & Faire retrouver cinq écrans de mémoire \\
1 h 30 & Service client (ch.~\ref{ch:responsable}) & Prendre, répondre, clore \\
1 h & Commandes et retraits & Refuser un retrait sans code \\
1 h & Catalogue et paliers & \textbf{L'exercice du palier fermé} \\
1 h & Droits (ch.~\ref{ch:admin}) & Créer un profil minimal, constater le menu \\
30 min & Journal (ch.~\ref{ch:superadmin}) & Modifier puis retrouver le geste \\
30 min & Checklists & Validation individuelle \\
\bottomrule
\end{tabular}
\end{center}

\section{Les erreurs à surveiller chez les stagiaires}

\begin{longtable}{@{}p{6cm} p{8.8cm}@{}}
\toprule
\bfseries Erreur fréquente & \bfseries Ce qu'il faut corriger \\
\midrule
\endhead
Confondre réclamation et retour
  & Faire ouvrir les deux écrans côte à côte \\[4pt]
Croire qu'un écran manquant est une panne
  & Rappeler que le menu ne montre que l'accessible \\[4pt]
Attendre que le paiement se confirme seul
  & Montrer le bouton qui redemande l'état, et dire pourquoi il existe \\[4pt]
Clôturer une conversation trop tôt
  & Faire constater que le client ne peut plus écrire \\[4pt]
Fermer le dernier palier de prix
  & Faire l'exercice du palier fermé \\[4pt]
Confirmer un retrait sans code
  & Rejouer le litige : deux paroles, aucune preuve \\[4pt]
Chercher dans le catalogue sans voir le filtre actif
  & Faire repérer la puce de filtre \\[4pt]
Donner tous les droits « pour simplifier »
  & Faire lister ce que la personne peut alors casser \\
\bottomrule
\end{longtable}

\section{Comment valider qu'un stagiaire est opérationnel}

Un stagiaire est opérationnel quand il coche \textbf{toute} la checklist de son
chapitre \emph{sans le manuel sous les yeux}, et qu'il sait répondre à ces trois
questions :

\begin{enumerate}[leftmargin=*,itemsep=4pt]
  \item « Un écran ne s'affiche pas dans votre menu. Que faites-vous ? »\\
        \textit{Attendu : ce n'est pas une panne ; demander le droit à un
        administrateur.}
  \item « Un client dit qu'il a écrit et vous ne voyez rien. Que vérifiez-vous
        d'abord ? »\\
        \textit{Attendu : la conversation est-elle prise en charge ?}
  \item « Que se passe-t-il si le dernier palier de prix est fermé ? »\\
        \textit{Attendu : au-delà, aucun prix n'existe et la vente est refusée.}
\end{enumerate}

% =============================================================================
\chapter*{Glossaire}
\addcontentsline{toc}{chapter}{Glossaire}
\markboth{Glossaire}{Glossaire}
% =============================================================================

\begin{description}[leftmargin=3.4cm,style=nextline,font=\bfseries]
  \item[Back-office] L'application interne, réservée aux responsables,
        administrateurs et super administrateurs.
  \item[Code de retrait] Le code unique qui prouve la remise de la marchandise.
        Seule preuve opposable.
  \item[Conversation] Un échange écrit entre un client et le service client.
        Trois états : en attente, en cours, fermée.
  \item[Déclinaison] Une version achetable d'un article. Porte le stock et le
        prix.
  \item[Expédition] Le voyage de colis sur un itinéraire.
  \item[Itinéraire] Le chemin suivi, avec ses points de transit.
  \item[Mobile Money] Le paiement par téléphone : MTN Mobile Money ou Orange
        Money.
  \item[Palier de prix] Une tranche de quantité associée à un prix unitaire.
  \item[Point de récupération] Le lieu où le client retire sa marchandise.
  \item[Profil] Un paquet de droits attribué à un responsable.
  \item[Proposition de prix] Un prix négocié dans une conversation. Elle expire.
  \item[Réclamation] Un problème signalé après la vente.
  \item[Retour] Une marchandise qui revient.
  \item[Service] Le regroupement métier d'un responsable.
  \item[Sous-total] Le montant des articles, \emph{hors} frais d'acheminement.
  \item[Total à payer] Le sous-total \emph{plus} les frais d'acheminement.
\end{description}

% =============================================================================
\chapter*{Conclusion}
\addcontentsline{toc}{chapter}{Conclusion}
\markboth{Conclusion}{Conclusion}
% =============================================================================

Vous avez maintenant tout ce qu'il faut pour travailler seul.

Si vous ne deviez retenir que trois choses de ce manuel :

\begin{enumerate}[leftmargin=*,itemsep=6pt]
  \item \textbf{La plateforme vous dit toujours pourquoi elle refuse.} Un
        bouton désactivé, une entrée absente du menu, un message d'erreur : ce
        ne sont pas des pannes, ce sont des réponses. Lisez-les avant de
        chercher plus loin.
  \item \textbf{Un droit, un code, un palier : chacun protège quelqu'un.} Le
        code de retrait protège le client \emph{et} l'agent. Le palier ouvert
        protège la vente. Le profil minimal protège l'entreprise.
  \item \textbf{Ce qui n'est pas tracé n'existe pas.} Le journal des actions est
        ce qui permet, des semaines plus tard, de répondre à « qui a fait ça ? »
        autrement que par des suppositions.
\end{enumerate}

\vfill
\begin{center}
\begin{tcolorbox}[width=0.8\linewidth,colback=garahFond,colframe=garahMarque,
  arc=5pt,boxrule=1.2pt]
\centering
{\Large\bfseries\color{garahEncre}%
 G\textcolor{garahMarque}{AR}\textcolor{garahAccent}{AH}}\\[3pt]
{\small\bfseries\color{garahSlogan}\MakeUppercase{Au-delà des frontières}}\\[8pt]
\footnotesize Douala $\rightarrow$ Bangui
\end{tcolorbox}
\end{center}
\vfill

\end{document}
